% Converted from Microsoft Word to LaTeX
% by Chikrii Softlab Word2TeX converter (version 3.0)
% Copyright (C) 1999-2003 Chikrii Softlab. All rights reserved.
% http://www.chikrii.com
% mailto: info@chikrii.com
% License: CSL#00004F

\documentclass{amsart}
\usepackage{latexsym}
\usepackage{amssymb}
\usepackage{amsmath}

\begin{document}
\begin{center}
\textbf{Two Body Propagator Approximation}
\end{center}

\begin{center}
\textbf{Using a Coherent States Approach}
\end{center}

\underline {\textbf{Introduction}}

Isolated quantum systems do not interact with their environment and are thus
unobservable. In order to determine their properties they must interact with
external probes such as electromagnetic fields of various frequencies. The
radiation field and the system under study, which we will consider to be a
molecular system, can be considered as a joint system that evolves under a
Hamiltonian that describes the behavior of both radiation and molecules.

The effect of the radiation on the molecular subsystem can be monitored by
changes in its associated density matrix and reduced density matrices. The
interaction of the electromagnetic field with molecules can be approximated
semi classically in a standard fashion by averaging over the photonic
degrees of freedom and only considering terms of first order in the electric
and magnetic fields. This leads to a reduced Hamiltonian for the molecular
subsystem that includes effective time dependent one body potentials
describing the interaction with the radiation field.

The behavior of a physical system, i.e. changes in its density matrices, is
fully described quantum mechanically by its associated propagator. A
powerful method that is often used to evaluate the full quantum propagator
between an initial and a final state is the path integral construction. Path
integrals can be formulated in terms of complex valued paths in spaces of
generalized coherent states. These spaces are complex differentiable
manifolds, that are also coset spaces of groups. The family of coherent
states chosen determines the specific geometric structure of the space of
the complex valued paths, \textbf{z}($t)$, via a metric function $C\left( {{\rm
{\bf z}},{\rm {\bf {\bar {z}}'}}} \right)$ . The time development of the
paths is described by a \textit{classical Hamiltonian }function $E\left( {{\rm {\bf \bar {z}}},{\rm {\bf
z}}} \right)$. The form and degree of these functions is determined by the
family of coherent states, and thus the group/coset pair. If the Lagrangian
associated with the Hamiltonian function is quadratic then the path integral
can be evaluated exactly, if this is not the case approximations must be
invoked.

If the external radiation fields are of low intensity and the molecular
system is in a stationary state, it is sufficient to consider the linear
response of the molecular state to the applied fields. This can be
characterized by a linear response function that can be calculated by using
first order perturbation theory. Noting that changes in the expectation
values of one particle observables are determined by changes in the one
particle reduced density matrix, one can describe the effects of the
radiation field, (approximated as an effective one particle interaction), on
the expectation values of any one particle observable in terms of retarded
two particle propagators.

Two particle propagators based on a given molecular initial state can be
expressed as a second order functional derivative of the propagator, that
contains interaction with arbitrary one particle external fields, with
respect to variations of these fields evaluated at an initial and a final
time. These times are chosen to bracket the time during which the molecular
system is actually interacting, before the initial time and after the final
time there is no interaction and the system is evolving "freely".

The quantum propagator of a N particle molecular system characterizes the
linear response of the system to external probes constituted of fields that
interact through M -body forces, where M $\leqslant $ N, as these are all
described by changes in the N - particle density matrix. Thus if one has
explicit analytic expressions for the full propagator of the molecular
system between arbitrary initial and final states one can deduce explicit
expressions for all \textit{reduced }propagators. If one is interested only to the response
of the system to one body external fields it is sufficient to consider two
particle propagators. This is true for both exact and approximate
propagators. As is well known exact solutions for quantum evolution
operators have only been determined for simple model systems, and are not
known for multi electron molecular systems that involve coulombic
interactions between electrons. Thus for practical applications the goal is
to develop systematic, cost effective, theoretically well founded and
nonarbitary approximation schemes, that lead to accurate enough results to
benefit experimentalists.

One way to obtain approximations to propagators is to approximate the path
integral by using the Stationary Phase Approximation (SPA) based on a
specific family of coherent states. This is equivalent to using the Time
Dependent Variation Principle (TDVP) based on the same family of states.
Approximate stationary states of the molecular system then correspond to the
special stationary phase "paths" that are critical points of the
corresponding \textit{classical }Hamiltonian. The linear response of these approximate
stationary states to external perturbing fields would then involve the full
\textit{classical }Hamiltonian, which again, however, usually makes the problem numerically
intractable.

In order to obtain computational feasible solutions one makes a first order
approximation to the linear response by considering a second order
approximation of the \textit{classical }Hamiltonian around a critical point. This construction
is analogous to including first order corrections to the SPA of the path
integral by expanding the \textit{classical }Hamiltonian in the Lagrangian to second order
around the classical path

The second order approximation to the \textit{classical }Hamiltonian clearly depends on the
coordinate system used, and hence on the family of coherent states being
utilized. If one is interested in the response to external fields that
interact through one body forces then these forces need to be expressible in
terms of the Lie algebra generators of the group associated with the
coherent states. For systems in which the number of particles is a good
quantum number and is conserved, i.e. a molecular system, the family of
coherent states should thus be based on the group SU(r), where r = 2s is the
size of the finite dimensional subspace one is approximating the full
Hilbert space by.

The Lie algebra of the group SU(r) has three Cartan decomposition into
\[
\begin{array}{l}
 su(2s)=so(2s,R)\oplus so(2s,R)^\bot \\
 su(2s)=su(N)\oplus su(2s-N)^\bot \\
 su(2s)=usp(2s)\oplus usp(2s)^\bot \\
 \end{array}
\]
\begin{center}
(1)
\end{center}

where the orthogonality is with respect to the killing form of \textit{su(2s)}. These
decomposition define three distinct coset decompositions of SU(r)
\[
\begin{array}{l}
 \raise0.7ex\hbox{${\mbox{SU(r)}}$} \!\mathord{\left/ {\vphantom
{{\mbox{SU(r)}} {\mbox{SU(N)}\otimes
\mbox{SU(r-N)}}}}\right.\kern-\nulldelimiterspace}\!\lower0.7ex\hbox{${\mbox{SU(N)}\otimes
\mbox{SU(r-N)}}$}\mbox{ } \\
 \raise0.7ex\hbox{${\mbox{SU(r)}}$} \!\mathord{\left/ {\vphantom
{{\mbox{SU(r)}}
{\mbox{USp(2s)}}}}\right.\kern-\nulldelimiterspace}\!\lower0.7ex\hbox{${\mbox{USp(2s)}}$}
\\
 \raise0.7ex\hbox{${\mbox{SU(r)}}$} \!\mathord{\left/ {\vphantom
{{\mbox{SU(r)}}
{\mbox{SO(r)}}}}\right.\kern-\nulldelimiterspace}\!\lower0.7ex\hbox{${\mbox{SO(r)}}$}
\\
 \end{array}
\]
\begin{center}
(2)
\end{center}

and thus three families of coherent states. The
$\raise0.7ex\hbox{${\mbox{SU(r)}}$} \!\mathord{\left/ {\vphantom
{{\mbox{SU(r)}} {\mbox{SU(N)}\otimes
\mbox{SU(r-N)}}}}\right.\kern-\nulldelimiterspace}\!\lower0.7ex\hbox{${\mbox{SU(N)}\otimes
\mbox{SU(r-N)}}$}\mbox{ }$- coherent states give the well known Thouless
parameterization of Independent Particle (IP) States and lead to the
standard Random Phase Approximation (RPA) for the linear response function.
The $\raise0.7ex\hbox{${\mbox{SU(r)}}$} \!\mathord{\left/ {\vphantom
{{\mbox{SU(r)}}
{\mbox{USp(2s)}}}}\right.\kern-\nulldelimiterspace}\!\lower0.7ex\hbox{${\mbox{USp(2s)}}$}$-
coherent states can be shown to be a parameterization of Antisymmetrized
Geminal Power (AGP) states and produce a generalized RPA for the linear
response function. At the present time the form of
$\raise0.7ex\hbox{${\mbox{SU(r)}}$} \!\mathord{\left/ {\vphantom
{{\mbox{SU(r)}}
{\mbox{SO(r)}}}}\right.\kern-\nulldelimiterspace}\!\lower0.7ex\hbox{${\mbox{SO(r)}}$}$-
coherent states are not known.

Coset decompositions of Lie groups when used in conjunction with Parabolic
Induction is a powerful method in the construction and characterization of
irreducible group representations. Representations are induced from the
action of the group on the coset space to the action of the group on
functions defined on the coset space. Points of the coset space corresponds
to coherent states, thus the induced representations act on functions of
coherent states. Hilbert spaces of these functions produce a holomorphic
representation of quantum states and thus of quantum mechanics. Due to the
coset structure and Lie group properties of the coherent state manifold,
quantum mechanical observables can be represented as differential operators
acting on functions defined on this manifold. By using such differential
operator representations one can obtain expressions for matrix elements of
observables in a coherent state \textit{overcomplete }basis and thus for expectation values for
observables with respect to any state. These expressions are of great
importance in practical applications.

The coset spaces are Homogeneous spaces that have the structure of Symmetric
spaces. The method of Parabolic Induction is a fundamental and powerful tool
used in the study of Noncommutative Harmonic Analysis on Symmetric spaces
and in the classification and properties of various classes of differential
operators that act on these spaces. The exact and approximate time evolution
and response of quantum states to external probes can be formulated within
these mathematical structures. Thus one has a complimentary between the
abstract classification, analysis and properties of these mathematical
structures on one hand and the quantitative description of concrete physical
processes on the other. The mathematical framework leads one to develop well
founded, rigorously expressed and improved approximation schemes, while the
physical constraints and properties imposed by real world applications
illuminates and focuses attention on specific relationships in the
mathematical structure that would not otherwise be apparent.

\underline {\textbf{Background}}

During the 1980's we developed approximation methods for two particle
propagators based on the properties of Antisymmtrized Geminal Power (AGP)
states. These properties were based on:

(a) Excited states, $\left\{ {\left| {\Psi _n } \right\rangle } \right\}$,
of the molecular system being well described by single particle excitations
$\left\{ {Q_n^\dag } \right\}$ from an energy optimized AGP state i.e.
\[
\left| {\Psi _n } \right\rangle =Q_n^\dag \left| {AGP} \right\rangle
=\sum\limits_{ij} {Q_{nij} a_i^\dag a_j^ \left| {AGP} \right\rangle }
\]
(b) a vacuum condition being satisfied
\[
Q_n^ \left| {AGP} \right\rangle =0
\]
(c) the off diagonal blocks of the energy Hessian at stationary points being
small in norm i.e.
\[
\begin{array}{c}
 E\left( {\zeta ,\bar {\zeta }} \right)=\frac{\left\langle {AGP\left( \zeta
\right)} \mathrel{\left| {\vphantom {{AGP\left( \zeta \right)} {H\mbox{
}AGP\left( \zeta \right)}}} \right. \kern-\nulldelimiterspace} {H\mbox{
}AGP\left( \zeta \right)} \right\rangle }{\left\langle {AGP\left( \zeta
\right)} \mathrel{\left| {\vphantom {{AGP\left( \zeta \right)} {\mbox{
}AGP\left( \zeta \right)}}} \right. \kern-\nulldelimiterspace} {\mbox{
}AGP\left( \zeta \right)} \right\rangle } \\
 {\rm {\bf D}}E\left( {\zeta _0 ,\bar {\zeta }_0 } \right)={\rm {\bf
0}}\mbox{; }{\rm {\bf D}}^{\rm {\bf 2}}E\left( {\zeta _0 ,\bar {\zeta }_0 }
\right)=\left( {{\begin{array}{*{20}c}
 {\frac{\partial ^2E\left( {\zeta _0 ,\bar {\zeta }_0 } \right)}{\partial
\zeta \partial \bar {\zeta }}} \hfill & {\frac{\partial ^2E\left( {\zeta _0
,\bar {\zeta }_0 } \right)}{\partial \zeta \partial \zeta }} \hfill \\
 {\frac{\partial ^2E\left( {\zeta _0 ,\bar {\zeta }_0 } \right)}{\partial
\bar {\zeta }\partial \bar {\zeta }}} \hfill & {\frac{\partial ^2E\left(
{\zeta _0 ,\bar {\zeta }_0 } \right)}{\partial \bar {\zeta }\partial \zeta
}} \hfill \\
\end{array} }} \right) \\
 \left\| {\frac{\partial ^2E\left( {\zeta _0 ,\bar {\zeta }_0 }
\right)}{\partial \zeta \partial \zeta }} \right\|<\varepsilon ;\mbox{
}\left\| {\frac{\partial ^2E\left( {\zeta _0 ,\bar {\zeta }_0 }
\right)}{\partial \bar {\zeta }\partial \bar {\zeta }}} \right\|<\varepsilon
\\
 \end{array}
\]
It was found that the approximation was improved by the inclusion of some
special many particle excitation operators that were associated with the
number operators $\left\{ {a_i^\dag a_i^ } \right\}$\underline {\textbf{,}}
however, the justification of including such operators was not well
understood. In the case of the exact propagator, for which $\left\{
{Q_n^\dag } \right\}$ are in general multiparticle operators, all of the
properties (a), (b) and (c) are satisfied with the off diagonal Hessian
block being identically zero i.e. ' =0.

The numerical results obtained from the application of this approximation
scheme agreed well with experiment, but many questions about the theoretical
basis of the approximation remained. Recently, however, this approximation
has been placed within the context of coherent states and the group
structure of many particle fermion state space. Within this context the
original approximation (including the modifications to improve it) can be
rigorously derived, and its relation to many other mathematical and physical
properties fully displayed. This new formulation also produces more
efficient and flexible numerical implementations.

\underline {\textbf{Proposed Work}}

We propose to:

1. Develop numerical procedures to find critical points of the generalized
classical equations of motion on the nonlinear phase space constituted by
$\raise0.7ex\hbox{${\mbox{SU(r)}}$} \!\mathord{\left/ {\vphantom
{{\mbox{SU(r)}}
{\mbox{USp(2s)}}}}\right.\kern-\nulldelimiterspace}\!\lower0.7ex\hbox{${\mbox{USp(2s)}}$}$-
coherent states. This corresponds to finding energy optimized AGP's. These
numerical procedures will be based on the theoretical work published in
"Coherent State Approach to Electron nuclear Dynamics with an
Antisymmetrized Geminal Power State". Which describes iterative schemes
based on nonorthogonal atomic orbital bases anchored to nuclear positions.
These schemes are similar to but more general than iterative Hartree-Fock
self-consistent field methods. The $\raise0.7ex\hbox{${\mbox{SU(r)}}$}
\!\mathord{\left/ {\vphantom {{\mbox{SU(r)}}
{\mbox{USp(2s)}}}}\right.\kern-\nulldelimiterspace}\!\lower0.7ex\hbox{${\mbox{USp(2s)}}$}$-
coherent states depend on nuclear positions through their atomic orbital
dependence, thus allowing the critical point search to simultaneously
directly find electronic energy and geometry optimized AGP states, as well
as transition states.

2. Analyze the foliation structure of the
$\raise0.7ex\hbox{${\mbox{SU(r)}}$} \!\mathord{\left/ {\vphantom
{{\mbox{SU(r)}}
{\mbox{USp(2s)}}}}\right.\kern-\nulldelimiterspace}\!\lower0.7ex\hbox{${\mbox{USp(2s)}}$}$-
coherent state manifold, and its implications for finding critical points.
This involves a study of the analytic properties of the metric function
$C\left( {{\rm {\bf z}},{\rm {\bf \bar {z}}}} \right)$.

3. Apply the method that is developed to practical problems of importance in
material science, physics and chemistry, in the study of the properties of
new materials and the clarification and understanding of experimental
spectra, characterizing new molecular products.

This will be accomplished by using energy optimized AGP states as
approximate stationary states in generalized RPA approximations to the
linear response function=retarded propagator and its Fourier-Laplace
transform. As the poles of the retarded propagator correspond to excitation
energies/frequencies, and the residues to oscillator strengths, approximate
UV, Visible, Infrared and microwave spectra can be constructed.

4. Study the representations of SL(r) and SU(r) induced by parabolic
induction based on the coset $\raise0.7ex\hbox{${\mbox{SU(r)}}$}
\!\mathord{\left/ {\vphantom {{\mbox{SU(r)}}
{\mbox{USp(2s)}}}}\right.\kern-\nulldelimiterspace}\!\lower0.7ex\hbox{${\mbox{USp(2s)}}$}$-
space and the physical meaning of the invariant differential operators that
act on holomorphic functions defined on this space.

In the following we will describe the technical details pertinent to this
proposal. Their presentation will be organized in the proceeding manner:

1. Definitions.

2. Path Integrals.

1. Coherent State Formulation.

2. Stationary Phase Approximation.

3. Coherent States.

1. Definitions and Construction.

2. Group Factorization.

3. Differential Structure of Coset Space.

4. Linear Response.

1. Exact Linear Response.

2. Linear Response in a Coherent State Setting.

3. Exact Linear Response in Coherent State Setting.

4. Linear Response Described by $\raise0.7ex\hbox{${\mbox{SU}\left( r
\right)}$} \!\mathord{\left/ {\vphantom {{\mbox{SU}\left( r \right)}
{\mbox{USp}\left( r
\right)}}}\right.\kern-\nulldelimiterspace}\!\lower0.7ex\hbox{${\mbox{USp}\left(
r \right)}$}$-Coherent States\underline {\textbf{.}}

5. Functions defined on the Coset Space.

1. Space of Holomorphic Quantum Wave Functions.

2. Representation of Observables as Differential Operators.

3. Induced Representations.

\newpage
\underline {\textbf{1. 0 Definitions.}}

We will base all expressions on the standard second quantized formulation
electronic quantum mechanics, i.e.
\[
\begin{array}{c}
 H_\mbox{F} =\mbox{ }\mathop \oplus \limits_{N=0}^{N=r} H^N\mbox{:
Fermi-Fock Space} \\
 H^N=\mbox{ }\Lambda H^1\mbox{ : N-fold Antisymmetric Tensor Product} \\
 H^1=\mbox{ linear span }\left\{ {\varphi _\mbox{i} ;1\leqslant i\leqslant
r} \right\};\mbox{ }\varphi _\mbox{i} \mbox{ are one particle functions of
position and spin} \\
 B(H_\mbox{F} )\mbox{ }=\mbox{ Bounded Operators acting in }H_\mbox{F} \\
 \left\{ {a\left( \varphi \right)^\dag ,a\left( \psi \right)}
\right\}=\left\langle {\varphi ,\psi } \right\rangle {\rm {\bf I}}\mbox{
Anticommutation Relationships} \\
 B\left( {H_\mbox{F} } \right)=\mbox{ linear Span }\left\{ {a_{i_1 }^{^\dag
} \cdots a_{i_p }^{^\dag } a_{i_1 }^ \cdots a_{i_q }^{^} } \right\};\mbox{
}1\leqslant i_1 <\cdots <i_p \leqslant r;\mbox{ }1\leqslant i_1 <\cdots <i_q
\leqslant r;\mbox{ 1}\leqslant p,q\leqslant r \\
 a_i^\dag =a\left( {\varphi _i } \right)^{^\dag };\mbox{ }\left\langle
{\varphi _i ,\varphi _j } \right\rangle =\delta _{ij} \\
 \end{array}
\]
\[
\begin{array}{l}
 \mbox{ The generators }\left\{ {a_i^\dag ,a_i^ ,a_i^\dag a_j
-\frac{1}{2}\delta _{ij} I,a_i^\dag a_j^\dag ,a_i a_j ;1\leqslant
i,j\leqslant r} \right\}\mbox{ satisfy the commutation } \\
 \mbox{relationships of the lie algebra }so\mbox{(2r+1), and generate a
representation } \\
 \mbox{of the universal enveloping algebra of }so\mbox{(2r+1) in
}B\mbox{(}H_F \mbox{), which is in fact all } \\
 \mbox{of }B\mbox{(}H_F \mbox{). The subset of generators }\left\{ {a_i^\dag
a_j -\frac{1}{2}\delta _{ij} I,a_i^\dag a_j^\dag ,a_i a_j ;1\leqslant
i,j\leqslant r} \right\}\mbox{ satisfy the } \\
 \mbox{commutation relationships of the lie algebra }so\mbox{(2r), and
generate a } \\
 \mbox{representation the universal enveloping of }so\mbox{(2r) in
}B\mbox{(}H_F \mbox{). The generators } \\
 \left\{ {a_i^\dag a_j -\frac{1}{2}\delta _{ij} I;1\leqslant i,j\leqslant r}
\right\}\mbox{ satisfy the commutation relationships of the lie algebra } \\
 su\mbox{(r), and generate a representation the universal enveloping of
}su\mbox{(r) in }B\mbox{(}H_F \mbox{). } \\
 \mbox{ The representation of }so\mbox{(2r+1) is irreducible, while the
representations of } \\
 so\mbox{(2r) and }su\mbox{(r) are reducible. Irreducible representations of
}so\mbox{(2r) are } \\
 \mbox{obtained by restricting to the subspaces of even or odd states in
}H_F \mbox{, while} \\
 \mbox{irreducible representations of }su\mbox{(r) are obtained by
restricting to the } \\
 \mbox{subspaces of states of fixed integer particle number. } \\
 \end{array}
\]
\newpage
\underline {\textbf{2. Path Integrals}}

\underline {\textbf{2.1. Coherent State Formulation .}}
\[
\begin{array}{c}
 \mbox{ Solutions of TDSE are of the form }\left| {\Psi \left( t \right)}
\right\rangle =e^{-\frac{i}{\hbar }Ht}\left| {\Psi \left( 0 \right)}
\right\rangle \mbox{. Introducing the resolution of} \\
 \mbox{ the identity }I=\int {\left| {\rm {\bf z}} \right\rangle
\left\langle {\rm {\bf z}} \right|} \mbox{ }d\mu \left( {\rm {\bf z}}
\right)\mbox{ the evolution operator matrix elements }\equiv
\mbox{propagator } \\
 \mbox{ kernel in this representation are } \\
 \mbox{ }P({\rm {\bf {z}'}},{\rm {\bf z}};t)=\left\langle {{\rm {\bf {z}'}}}
\mathrel{\left| {\vphantom {{{\rm {\bf {z}'}}} {e^{-\frac{i}{\hbar }Ht}{\rm
{\bf z}}}}} \right. \kern-\nulldelimiterspace} {e^{-\frac{i}{\hbar }Ht}{\rm
{\bf z}}} \right\rangle \mbox{ and }\left| {\Psi \left( t \right)}
\right\rangle =\int {\left| {{\rm {\bf {z}'}}} \right\rangle P({\rm {\bf
{z}'}},{\rm {\bf z}};t)\left\langle {\rm {\bf z}} \right|d\mu \left( {{\rm
{\bf {z}'}}} \right)} d\mu \left( {\rm {\bf z}} \right)\mbox{ }\left| {\Psi
\left( 0 \right)} \right\rangle \\
 \end{array}
\]
\[
\begin{array}{l}
 \mbox{Using} \mbox{the} \mbox{coherent state construction} \mbox{the}
\mbox{Propagator} \mbox{kernel} becomes \mbox{an } \\
 \mbox{integral} \mbox{over} \mbox{paths with} \mbox{fixed} \mbox{end}
\mbox{points }{\rm {\bf z}}\left( 0 \right)= {\rm {\bf z}}_i \mbox{ and
}{\rm {\bf z}}\left( t \right)= {\rm {\bf z}}_f \\
 P \left( {{\rm {\bf z}}_f ,{\rm {\bf z}}_i ,t} \right)\approx \int_{{\rm
{\bf z}}_i }^{{\rm {\bf z}}_f } {e^{\left\{ {\frac{i}{\hbar }\int_0^t
{\left\{ {\left. {\frac{i\hbar }{2}\ln \left\langle {{\rm {\bf z}}\left(
{{{t}'}'} \right)} \mathrel{\left| {\vphantom {{{\rm {\bf z}}\left(
{{{t}'}'} \right)}
{\overset{\lower0.5em\hbox{$\smash{\scriptscriptstyle\leftrightarrow}$}}\to
{{\partial }} _t {\rm {\bf z}}\left( {{t}'} \right)}}} \right.
\kern-\nulldelimiterspace}
{\overset{\lower0.5em\hbox{$\smash{\scriptscriptstyle\leftrightarrow}$}}\to
{{\partial }} _t {\rm {\bf z}}\left( {{t}'} \right)} \right\rangle }
\right|_{{{t}'}'={t}'} -E\left( {{\rm {\bf z}}\left( {{t}'} \right)}
\right)} \right\}d\mbox{ }{t}'} } \right\}}d\mu ({\rm {\bf z}})} \mbox{, }
\\
 \mbox{where }E\left( {{\rm {\bf z}}\left( {{t}'} \right)}
\right)=\frac{\left\langle {{\rm {\bf z}}\left( {{t}'} \right)}
\mathrel{\left| {\vphantom {{{\rm {\bf z}}\left( {{t}'} \right)} {H{\rm {\bf
z}}\left( {{t}'} \right)}}} \right. \kern-\nulldelimiterspace} {H{\rm {\bf
z}}\left( {{t}'} \right)} \right\rangle }{\left\langle {{\rm {\bf z}}\left(
{{t}'} \right)} \mathrel{\left| {\vphantom {{{\rm {\bf z}}\left( {{t}'}
\right)} {{\rm {\bf z}}\left( {{t}'} \right)}}} \right.
\kern-\nulldelimiterspace} {{\rm {\bf z}}\left( {{t}'} \right)}
\right\rangle } \\
 \end{array}
\]
The path integral Action and Largrangian are defined as
\[
\begin{array}{c}
 P\left( {{\rm {\bf z}}_f ,{\rm {\bf z}}_i ,t} \right)=\int_{{\rm {\bf z}}_i
}^{{\rm {\bf z}}_f } {e^{\frac{i}{\hbar }A\left( {{\rm {\bf z}}(t)}
\right)}d\mu ({\rm {\bf z}})} \mbox{ } \\
 A\left( {{\rm {\bf z}}(t)} \right)=\int_0^t {L\left( {{\rm {\bf z}},{\rm
{\bf \dot {z}}},{t}'} \right)d{t}'} \\
 L\left( {{\rm {\bf z}},{\rm {\bf \dot {z}}},{t}'} \right)=\left.
{\frac{i\hbar }{2}\ln \left\langle {{\rm {\bf z}}\left( {{{t}'}'} \right)}
\mathrel{\left| {\vphantom {{{\rm {\bf z}}\left( {{{t}'}'} \right)}
{\overset{\lower0.5em\hbox{$\smash{\scriptscriptstyle\leftrightarrow}$}}\to
{{\partial }} _t {\rm {\bf z}}\left( {{t}'} \right)}}} \right.
\kern-\nulldelimiterspace}
{\overset{\lower0.5em\hbox{$\smash{\scriptscriptstyle\leftrightarrow}$}}\to
{{\partial }} _t {\rm {\bf z}}\left( {{t}'} \right)} \right\rangle }
\right|_{{{t}'}'={t}'} -E\left( {{\rm {\bf z}}\left( {{t}'} \right)} \right)
\\
 \end{array}
\]
\underline {\textbf{2.2 Stationary Phase Approximation}}
\[
\begin{array}{l}
 \mbox{The stationary phase approximation }\delta A\left( {{\rm {\bf z}}_c }
\right)=0;\mbox{ with fixed end points }{\rm {\bf z}}_c (t)={\rm {\bf z}}_f
; \\
 \mbox{ }{\rm {\bf z}}_c (0)={\rm {\bf z}}_i \mbox{ produces the generalized
classical equations } \\
 \mbox{ }i\hbar \sum {C_{ij} \frac{dz_j }{dt}} =\frac{\partial E}{\partial
\bar {z}_i };\mbox{ -}i\hbar \sum {\bar {C}_{ij} \frac{d\bar {z}_j }{dt}}
=\frac{\partial E}{\partial z_i } \\
 \mbox{over the phase space formed by the manifold of states }\left\{
{\left| {\rm {\bf z}} \right\rangle } \right\}\mbox{. } \\
 \end{array}
\]
\[
\begin{array}{c}
 \mbox{Applying the Stationary Phase Approximation one obtains } \\
 \int_{{\rm {\bf z}}_i }^{{\rm {\bf z}}_f } {e^{\frac{i}{\hbar }A\left(
{{\rm {\bf z}}(t)} \right)}d\mu ({\rm {\bf z}})} \approx \sum\limits_c {K_c
\left| {\det \left( {\frac{i}{2\pi \hbar }\frac{\partial ^2A\left( {{\rm
{\bf z}}_c } \right)}{\partial {\rm {\bf x}}\partial {\rm {\bf x}}}}
\right)} \right|^{-\frac{1}{2}}e^{\frac{i}{\hbar }A\left( {{\rm {\bf z}}_c
(t)} \right)}} \\
 \mbox{where }{\rm {\bf x}}\mbox{=}{\rm {\bf z}}\mbox{ or }{\rm {\bf \bar
{z}}},\mbox{ with the conditions }\delta A\left( {\rm {\bf z}}
\right)=0\Leftrightarrow \delta L\left( {\rm {\bf z}} \right)=0;\mbox{ }{\rm
{\bf z}}_c (t)={\rm {\bf z}}_f ;\mbox{ }{\rm {\bf z}}_c (0)={\rm {\bf z}}_i
\\
 \mbox{and the paths }{\rm {\bf z}}_c \left( t \right)\mbox{ are solutions
to generalized classical equations of motion. } \\
 \end{array}
\begin{array}{c}
 \mbox{ If the Lagragian is Quadratic the "classical paths" }\left\{ {{\rm
{\bf z}}_c (t)} \right\}\mbox{ produce an } \\
 \mbox{exact propagator given by } \\
 P\left( {{\rm {\bf z}}_f ,{\rm {\bf z}}_i ,t} \right)=\sum\limits_c {\sqrt
{\det \left\{ {\frac{i}{h}\frac{\partial ^2A({\rm {\bf z}}_c (t))}{\partial
{\rm {\bf z}}_f \partial {\rm {\bf z}}_i }} \right\}} } e^{\frac{i}{\hbar
}A\left( {{\rm {\bf z}}_{\rm {\bf c}} \left( t \right)} \right)} \\
 \end{array}
\]
\[
\begin{array}{c}
 \mbox{So for stationary phase approximation or quadratic Lagragians} \\
 \left\langle {{\rm {\bf {z}'}}} \mathrel{\left| {\vphantom {{{\rm {\bf
{z}'}}} {\Psi \left( t \right)}}} \right. \kern-\nulldelimiterspace} {\Psi
\left( t \right)} \right\rangle \approx \int {\sum\limits_c {\Omega \left(
{{\rm {\bf z}}_c \left( t \right)} \right)e^{\frac{i}{\hbar }A_{{\rm {\bf
{z}'z}}} \left( {{\rm {\bf z}}_c (t)} \right)}} } \left\langle {{\rm {\bf
z}}} \mathrel{\left| {\vphantom {{\rm {\bf z}} {\Psi \left( 0 \right)}}}
\right. \kern-\nulldelimiterspace} {\Psi \left( 0 \right)} \right\rangle
\mbox{ }d\mu \left( {\rm {\bf z}} \right) \\
 \Updownarrow \\
 \left\langle {{\rm {\bf z}}\left( t \right)} \mathrel{\left| {\vphantom
{{{\rm {\bf z}}\left( t \right)} \Psi }} \right. \kern-\nulldelimiterspace}
{\Psi } \right\rangle =\int {\sum\limits_c {\Omega \left( {{\rm {\bf z}}_c
\left( t \right)} \right)e^{\frac{i}{\hbar }A_{{\rm {\bf {z}'z}}} \left(
{{\rm {\bf z}}_c (t)} \right)}} } \left\langle {{\rm {\bf z}}\left( 0
\right)} \mathrel{\left| {\vphantom {{{\rm {\bf z}}\left( 0 \right)} \Psi }}
\right. \kern-\nulldelimiterspace} {\Psi } \right\rangle \mbox{ }d\mu \left(
{\rm {\bf z}} \right) \\
 \Updownarrow \\
 \Psi \left( {{\rm {\bf z}}\left( t \right)} \right)=\int {\sum\limits_c
{\Omega \left( {{\rm {\bf z}}_c \left( t \right)} \right)e^{\frac{i}{\hbar
}A_{{\rm {\bf {z}'z}}} \left( {{\rm {\bf z}}_c (t)} \right)}} } \Psi \left(
{{\rm {\bf z}}\left( t \right)} \right)\mbox{ }d\mu \left( {\rm {\bf z}}
\right) \\
 \end{array}
\]
The SPA approximation can be improved by using a quadratic approximation to
the classical Hamiltonian around stationary paths, produced by truncating
the Taylor series expansion about such points at second order.

\underline {\textbf{3. Coherent States}}

\underline {\textbf{3.1 Definition and Construction}}

In a quantum mechanical context the following prescription is the one
essentially followed in the construction of group coherent states.
\[
\begin{array}{l}
 \mbox{ The coherent state construction is based on finding a subspace of
the hilbert space of } \\
 \mbox{states }H, \mbox{that carries a unitary irrep, G}_H \mbox{, of a
compact group G. The subspace can be } \\
 \mbox{all of }H.\mbox{ } \\
 \mbox{ The groups we will be considering are SU}\left( {H^\mbox{N}\mbox{ }}
\right)\equiv \mbox{SU}\left( {\left( {{\begin{array}{*{20}c}
 \mbox{r} \hfill \\
 \mbox{N} \hfill \\
\end{array} }} \right)\mbox{ }} \right)\mbox{ and SU}\left( {H^1\mbox{ }}
\right)\equiv \mbox{SU}\left( \mbox{r} \right),\mbox{ } \\
 \mbox{the groups associated with N-electron state space and 1-electron
state space respectively. } \\
 \mbox{The irreducibel representations of SU}\left( {H^\mbox{N}\mbox{ }}
\right)\mbox{ and SU}\left( {H^1\mbox{ }} \right)\mbox{ are those carried by
}H^\mbox{N}\mbox{. } \\
 \mbox{ A reference} \mbox{state, }\Phi \in H, \mbox{determines a maximally
compact "isotropy" group, } \\
 \mbox{H}_\Phi \mbox{ of }\Phi \mbox{ that is contained in G}_H \subseteq
\mbox{SU}\left( {H\mbox{ }} \right)\mbox{, where H}_\Phi \mbox{=}\left\{
{\mbox{g}\in \mbox{G}_H \mbox{; g :}\Phi \to \Phi } \right\}\mbox{. G}_H
\mbox{ is then } \\
 \mbox{factored through the left coset so that G}_H =\left( {{\mbox{G}_H }
\mathord{\left/ {\vphantom {{\mbox{G}_H } {\mbox{H}_V }}} \right.
\kern-\nulldelimiterspace} {\mbox{H}_V }} \right)\mbox{H}_V .\mbox{ } \\
 \mbox{ The reference state we will be focusing on in the proposed work is
an } \\
 \mbox{Antisymmetrized Geminal Power (AGP) state }\left|
{g^{\frac{\mbox{N}}{\mbox{2}}}} \right\rangle \in H^\mbox{N}\mbox{ given by
} \\
 \mbox{ }g^{\frac{\mbox{N}}{\mbox{2}}}=\underbrace {\mbox{g}\Lambda
\mbox{g}\cdots \mbox{g}\Lambda
\mbox{g}}_{\frac{\mbox{N}}{\mbox{2}}-\mbox{factors}};\mbox{ where
g=}\sum\limits_{1\leqslant i\leqslant s=\raise0.7ex\hbox{$r$}
\!\mathord{\left/ {\vphantom {r
2}}\right.\kern-\nulldelimiterspace}\!\lower0.7ex\hbox{$2$}} {\xi _\mbox{i}
} \left| {\psi _i \psi _{-i} } \right\rangle ;\mbox{ }\xi _\mbox{i} \in {\rm
{\bf R}}\mbox{ and }\left\{ {\psi _i ;\mbox{ -}\raise0.7ex\hbox{$r$}
\!\mathord{\left/ {\vphantom {r
2}}\right.\kern-\nulldelimiterspace}\!\lower0.7ex\hbox{$2$}\leqslant
i\leqslant \raise0.7ex\hbox{$r$} \!\mathord{\left/ {\vphantom {r
2}}\right.\kern-\nulldelimiterspace}\!\lower0.7ex\hbox{$2$}} \right\}\mbox{
is a conb for }H^\mbox{1} \\
 \mbox{and are the cannonical general spin orbitals of }g\mbox{ i.e. those
that bring it to cannonical } \\
 \mbox{form. This state can also be expressed as } \\
 \mbox{ }\left| {g^{\frac{\mbox{N}}{\mbox{2}}}} \right\rangle =P_\mbox{N}
e^{\sum\limits_{1\leqslant i\leqslant s=\raise0.7ex\hbox{$r$}
\!\mathord{\left/ {\vphantom {r
2}}\right.\kern-\nulldelimiterspace}\!\lower0.7ex\hbox{$2$}} {\xi _\mbox{i}
\mbox{b}_i^\dag \mbox{b}_{-i}^\dag } }\left| \phi \right\rangle \mbox{=
}P_\mbox{N} e^{\sum\limits_{1\leqslant i\leqslant s=\raise0.7ex\hbox{$r$}
\!\mathord{\left/ {\vphantom {r
2}}\right.\kern-\nulldelimiterspace}\!\lower0.7ex\hbox{$2$}} {g_{ij}
\mbox{a}_i^\dag \mbox{a}_j^\dag } }\left| \phi \right\rangle ;\mbox{
b}_i^\dag =\mbox{a}\left( {\psi _i } \right)_^\dag \\
 \mbox{where }P_\mbox{N} \mbox{ is the orthogonal projector onto
}H^\mbox{N}\mbox{ and }\left| \phi \right\rangle \mbox{ is the zero particle
vacuum state in }H_F .\mbox{ } \\
 \mbox{ } \\
 \end{array}
\]
\underline {\textbf{2.2. Group Factorization}}

The isotropy group for an AGP reference state is $\mbox{USp(}H^1\mbox{)}$
which is a maximally compact subgroup of $\mbox{ SU}\left( {H^1\mbox{ }}
\right)$ hence producing the coset
$\raise0.7ex\hbox{${\mbox{SU(}H^1\mbox{)}}$} \!\mathord{\left/ {\vphantom
{{\mbox{SU(}H^1\mbox{)}}
{\mbox{USp(}H^1\mbox{)}}}}\right.\kern-\nulldelimiterspace}\!\lower0.7ex\hbox{${\mbox{USp(}H^1\mbox{)}}$}$.
One can express coset representatives in factored form as
$\mbox{G}_{\mbox{+1}} \mbox{G}_\mbox{0} \mbox{G}_{\mbox{-1}} $ , where
$\left\{ {\mbox{G}_{\mbox{+1}} \mbox{,G}_\mbox{0} \mbox{,G}_{\mbox{-1}} }
\right\}$ are subgroups of $\mbox{ SL}\left( {H^1\mbox{ }} \right)$ and
$\mbox{G}_{\mbox{-1}} $ leaves the reference state invariant. Restricting
attention to the connected components of these subgroups, that contain
identity elements, one can express this factorization in terms of Lie
algebras as
\[
\mbox{exp}\left\{ {N_{+1} } \right\}\mbox{exp}\left\{ {N_0 }
\right\}\mbox{exp}\left\{ {N_{-1} } \right\}
\]
where $\left\{ {N_{+1} ,N_0 ,N_{-1} } \right\}$ are Lie subalgebras of
$sl\left( {H^1\mbox{ }} \right)$. Introducing coordinate systems $\left\{
{{\rm {\bf z}}_+ ,\lambda ,{\rm {\bf z}}_- ,\eta } \right\}$ with respect to
fixed bases of $\left\{ {N_{+1} ,N_0 ,N_{-1} ,h} \right\}$, where $\left\{ h
\right\}$ is the isotropy subalgebra, one can parameterize group elements
$g\in \mbox{SU}\left( {H^1\mbox{ }} \right)$ as
\[
\mbox{g}\left( {{\rm {\bf z}},\eta } \right)\mbox{=c}\left( {\rm {\bf z}}
\right)\mbox{h}\left( \eta \right)\mbox{=c}_\mbox{+} \left( {{\rm {\bf z}}_+
} \right)\mbox{c}_\mbox{0} \left( \lambda \right)\mbox{ c}_{\mbox{-1}}
\left( {{\rm {\bf z}}_- } \right)\mbox{h}\left( \eta \right)
\]
the non redundant coordinates $\left\{ {{\rm {\bf z}}\equiv ({\rm {\bf z}}_+
,\lambda )} \right\}$ describe a configuration space that parameterize the
set of $\left\{ {\left| {\rm {\bf z}} \right\rangle } \right\}\mbox{ of
}{\mbox{SU}\left( {H^1} \right)} \mathord{\left/ {\vphantom
{{\mbox{SU}\left( {H^1} \right)} {\mbox{USp}\left( {H^1} \right)}}} \right.
\kern-\nulldelimiterspace} {\mbox{USp}\left( {H^1} \right)}$ coherent states
and at the same time elements of the coset ${\mbox{SU}\left( {H^1} \right)}
\mathord{\left/ {\vphantom {{\mbox{SU}\left( {H^1} \right)}
{\mbox{USp}\left( {H^1} \right)}}} \right. \kern-\nulldelimiterspace}
{\mbox{USp}\left( {H^1} \right)}$by
\[
\mbox{g}\left| {g^{\tfrac{\mbox{N}}{\mbox{2}}}} \right\rangle
\mbox{=c}_\mbox{+} \left( {{\rm {\bf z}}_+ } \right)\mbox{c}_\mbox{0} \left(
\lambda \right)\mbox{c}_{\mbox{-1}} \left( {{\rm {\bf z}}_- }
\right)\mbox{h}\left( \eta \right)\left| {g^{\tfrac{\mbox{N}}{\mbox{2}}}}
\right\rangle =\mbox{c}_\mbox{+} \left( {{\rm {\bf z}}_+ }
\right)\mbox{c}_\mbox{0} \left( \lambda \right)\left|
{g^{\tfrac{\mbox{N}}{\mbox{2}}}} \right\rangle
\]
$\mbox{The set }\left\{ {\left| {\rm {\bf z}} \right\rangle
\mbox{=c}_\mbox{+} \left( {{\rm {\bf z}}_+ } \right)\mbox{c}_\mbox{0} \left(
\lambda \right)\left| {g^{\tfrac{\mbox{N}}{\mbox{2}}}} \right\rangle }
\right\}\mbox{ of }$coherent states forms a manifold $M$ (in general
nonlinear) in $H^{\left( {{\begin{array}{*{20}c}
 \mbox{r} \hfill \\
 \mbox{N} \hfill \\
\end{array} }} \right)} \quad \mbox{with a coordinate system }{\rm {\bf z}}$ .
This manifold is a representation of the coset space ${\mbox{SU}\left( {H^1}
\right)} \mathord{\left/ {\vphantom {{\mbox{SU}\left( {H^1} \right)}
{\mbox{USp}\left( {H^1} \right)}}} \right. \kern-\nulldelimiterspace}
{\mbox{USp}\left( {H^1} \right)}$

The bases one chooses are
\[
\begin{array}{c}
 N_{+1} =\left\{ {\xi _i a_j^\dag a_i +\sigma \left( i \right)\sigma \left(
j \right)\xi _j a_{-i}^\dag a_{-j} \mbox{; }-s\leqslant i<j\leqslant
s;\left| i \right|\ne \left| j \right|\mbox{; }\sigma \left( i
\right)=\mbox{sign}(i)} \right\} \\
 N_0 =\left\{ {a_i^\dag a_j -\frac{1}{2}\delta _{ij} I;\mbox{-}s\leqslant
i\leqslant s\mbox{; }\left| i \right|=\left| j \right|} \right\} \\
 N_{-1} =\left\{ {\xi _i a_i^\dag a_j -\sigma \left( i \right)\sigma \left(
j \right)\xi _j a_{-j}^\dag a_{-i} ;\mbox{ }-s\leqslant i<j\leqslant
s;\left| i \right|\ne \left| j \right|\mbox{ }} \right\} \\
 \eta =\left\{ {a_i^\dag a_j -\frac{1}{2}\delta _{ij} I;\mbox{-}s\leqslant
i\leqslant s\mbox{; }\left| i \right|=\left| j \right|} \right\} \\
 \end{array}
\]
and the factors in the group decomposition are
\[
\begin{array}{c}
 \mbox{c}_\mbox{+} \left( {{\rm {\bf z}}_+ }
\right)\mbox{=exp}\sum\limits_{\begin{array}{l}
 -s\leqslant i<j\leqslant s \\
 \left| i \right|\ne \left| j \right| \\
 \end{array}} {\mbox{z}_\mbox{+} _{ij} \left( {\xi _i a_j^\dag a_i +\sigma
\left( i \right)\sigma \left( j \right)\xi _j a_{-i}^\dag a_{-j} }
\right)=\mbox{exp}\sum\limits_{\begin{array}{l}
 -s\leqslant i<j\leqslant s \\
 \left| i \right|\ne \left| j \right| \\
 \end{array}} {\mbox{z}_\mbox{+} _{ij} q_{ij}^\dag } } \\
 \mbox{c}_{\mbox{-1}} \left( {{\rm {\bf z}}_{-1} }
\right)=\mbox{exp}-\sum\limits_{\begin{array}{l}
 -s\leqslant i<j\leqslant s \\
 \left| i \right|\ne \left| j \right| \\
 \end{array}} {\mbox{z}_{\mbox{-1}} _{ij} \left( {\xi _i a_i^\dag a_j
-\sigma \left( i \right)\sigma \left( j \right)\xi _j a_{-j}^\dag a_{-i} }
\right)} =\mbox{exp}\sum\limits_{\begin{array}{l}
 -s\leqslant i<j\leqslant s \\
 \left| i \right|\ne \left| j \right| \\
 \end{array}} {\mbox{z}_\mbox{-} _{ij} q_{ij}^ } \\
 \mbox{ } \\
 \mbox{c}_\mbox{0} \left( \lambda
\right)=\mbox{exp}\sum\limits_{\begin{array}{l}
 \mbox{-}s\leqslant i\leqslant s \\
 \left| i \right|=\left| j \right| \\
 \end{array}} {\left\{ {\lambda _{ij} \left( {a_i^\dag a_j
-\frac{1}{2}\delta _{ij} I} \right)}
\right\}=\mbox{exp}\sum\limits_{\begin{array}{l}
 \mbox{-}s\leqslant i\leqslant s \\
 \left| i \right|=\left| j \right| \\
 \end{array}} {\left\{ {\lambda _{ij} u_{ij}^ } \right\}} } ;\mbox{ }\lambda
_{ij} =\bar {\lambda }_{ji} \\
 \mbox{h}\left( \eta \right)=\mbox{exp}\sum\limits_{\begin{array}{c}
 -s\leqslant i,j\leqslant s \\
 \left| i \right|=\left| j \right| \\
 \end{array}} {\left\{ {\eta _{ij} \left( {a_i^\dag a_j -\frac{1}{2}\delta
_{ij} I} \right)} \right\}=\mbox{exp}\sum\limits_{\begin{array}{l}
 \mbox{-}s\leqslant i\leqslant s \\
 \left| i \right|=\left| j \right| \\
 \end{array}} {\left\{ {\eta _{ij} u_{ij}^ } \right\}} ;\mbox{ }} \eta _{ij}
=-\bar {\eta }_{ji} \\
 \end{array}
\]
Normalized coherent states are produced by the action of $\mbox{G}_H $ on
normalized reference states and thus by the action of the coset
representatives. Often however it is convenient to consider \textit{unnormalized }coherent states
by either dropping or modifying the coset representative $\mbox{c}_\mbox{0}
\left( \lambda \right)$ and explicitly considering the norm of the state.

\underline {\textbf{2.3. Differential Structure of the Coset Space}}

The phase space associated with the coset space $\mbox{G}_H \mbox{/H}_H
$inherits a coordinate system $\left\{ {{\rm {\bf z}},{\rm {\bf \bar {z}}}}
\right\}$ from the coset space and is a Kaehlerian manifold with a closed
symplectic form defined as
\[
\begin{array}{c}
 \omega \mbox{=-}\frac{i\hbar }{2}\sum {\left. {\left\{ {\frac{\partial
\mbox{ }^2F}{\partial \mbox{\bar {{z}'}}_\mbox{i} \partial \mbox{z}_\mbox{j}
}} \right\}} \right|} _{{\rm {\bf z}}={\rm {\bf {z}'}}} \mbox{d}{\rm {\bf
z}}\Lambda \mbox{d}{\rm {\bf \bar {z}}};\mbox{ } \\
 \mbox{which acts on tangent vectors }\left\{ X \right\}\mbox{ of this
manifold by: }\omega \left( {X_i ,X_j } \right)=-\omega \left( {X_j ,X_i }
\right)\mbox{. It is induced by } \\
 \mbox{the Kaehler potential }F\left( {{\rm {\bf {z}'}},{\rm {\bf z}}}
\right)\mbox{=}\ln \mbox{ }\left\langle {{\rm {\bf {z}'}}} \mathrel{\left|
{\vphantom {{{\rm {\bf {z}'}}} {\rm {\bf z}}}} \right.
\kern-\nulldelimiterspace} {{\rm {\bf z}}} \right\rangle ,\mbox{ and defines
a generalized Poisson Bracket by } \\
 \left\{ {f,g} \right\}=\omega \left( {X_f ,X_g } \right);\mbox{ where }X_f
=\left( {{\begin{array}{*{20}c}
 {\frac{\partial f}{\partial {\rm {\bf \bar {z}}}}} \hfill \\
 {\frac{\partial f}{\partial {\rm {\bf z}}}} \hfill \\
\end{array} }} \right) \\
 \mbox{ } \\
 \end{array}
\]
\[
\begin{array}{c}
 \mbox{The generalized classical equations of motion defined on this phase
space can be } \\
 \mbox{then expressed as } \\
 \frac{i}{\hbar }{\rm {\bf \dot {z}}}_i =\left\{ {{\rm {\bf z}}_i ,E}
\right\}\mbox{ }\And \mbox{ -}\frac{i}{\hbar }{\rm {\bf \dot {\bar {z}}}}_i
=\left\{ {{\rm {\bf \bar {z}}}_i ,E} \right\} \\
 \mbox{which can be written in matrix form as } \\
 \left( {{\begin{array}{*{20}c}
 {{\rm {\bf \dot {\bar {z}}}}} \hfill \\
 {\rm {\bf z}} \hfill \\
\end{array} }} \right)=\left( {{\begin{array}{*{20}c}
 {\rm {\bf 0}} \hfill & {\tfrac{i}{\hbar }{\rm {\bf \bar {C}}}^{-{\rm {\bf
1}}}} \hfill \\
 {-\tfrac{i}{\hbar }{\rm {\bf C}}^{-{\rm {\bf 1}}}} \hfill & {\rm {\bf 0}}
\hfill \\
\end{array} }} \right)\left( {{\begin{array}{*{20}c}
 {\frac{\partial E}{\partial {\rm {\bf \bar {z}}}}} \hfill \\
 {\frac{\partial E}{\partial {\rm {\bf z}}}} \hfill \\
\end{array} }} \right). \\
 \end{array}
\]
\[
\begin{array}{c}
 \mbox{The metric matrix }{\rm {\bf C}}=\left. {\frac{\partial ^2\ln
\left\langle {{\rm {\bf {z}'}}} \mathrel{\left| {\vphantom {{{\rm {\bf
{z}'}}} {\rm {\bf z}}}} \right. \kern-\nulldelimiterspace} {{\rm {\bf z}}}
\right\rangle }{\partial {\rm {\bf {\bar {z}}'}}\partial {\rm {\bf z}}}}
\right|_{{z}'=z} \mbox{ is given by } \\
 {\rm {\bf C}}=\frac{1}{S}\left\{ {\left. {\frac{\partial ^2\left\langle
{{\rm {\bf {z}'}}} \mathrel{\left| {\vphantom {{{\rm {\bf {z}'}}} {\rm {\bf
z}}}} \right. \kern-\nulldelimiterspace} {{\rm {\bf z}}} \right\rangle
}{\partial {\rm {\bf {\bar {z}}'}}\partial {\rm {\bf z}}}} \right|_{{z}'=z}
-\frac{1}{S}\left. {\frac{\partial \left\langle {{\rm {\bf {z}'}}}
\mathrel{\left| {\vphantom {{{\rm {\bf {z}'}}} {\rm {\bf z}}}} \right.
\kern-\nulldelimiterspace} {{\rm {\bf z}}} \right\rangle }{\partial {\rm
{\bf {\bar {z}}'}}}} \right|_{{z}'=z} \left. {\frac{\partial \left\langle
{{\rm {\bf {z}'}}} \mathrel{\left| {\vphantom {{{\rm {\bf {z}'}}} {\rm {\bf
z}}}} \right. \kern-\nulldelimiterspace} {{\rm {\bf z}}} \right\rangle
}{\partial {\rm {\bf z}}}} \right|_{{z}'=z} } \right\}\equiv
\frac{1}{S}\left\{ {\left. {\frac{\partial ^2S\left( {{\rm {\bf z\bar
{{z}'}}}} \right)}{\partial {\rm {\bf {\bar {z}}'}}\partial {\rm {\bf z}}}}
\right|_{{z}'=z} -\frac{1}{S}\left. {\frac{\partial S\left( {{\rm {\bf z\bar
{{z}'}}}} \right)}{\partial {\rm {\bf {\bar {z}}'}}}} \right|_{{z}'=z}
\left. {\frac{\partial S\left( {{\rm {\bf z\bar {{z}'}}}} \right)}{\partial
{\rm {\bf z}}}} \right|_{{z}'=z} } \right\} \\
 \mbox{while the gradient of the energy function }E\left( {{\rm {\bf
z}},{\rm {\bf \bar {z}}}} \right)\mbox{ has the form } \\
 {\rm {\bf D}}E\left( {{\rm {\bf z}},{\rm {\bf \bar {z}}}}
\right)=\frac{\partial E\left( {{\rm {\bf z}},{\rm {\bf \bar {z}}}}
\right)}{\partial {\rm {\bf x}}}=\frac{1}{S}\left\{ {\frac{\partial H\left(
{{\rm {\bf z}},{\rm {\bf \bar {z}}}} \right)}{\partial {\rm {\bf
x}}}-E\frac{\partial S\left( {{\rm {\bf z}},{\rm {\bf \bar {z}}}}
\right)}{\partial {\rm {\bf x}}}} \right\};\mbox{ where }{\rm {\bf
x}}=\left( {{\rm {\bf z}},{\rm {\bf \bar {z}}}} \right) \\
 \mbox{at critical points }\left( {{\rm {\bf z}}_0 ,{\rm {\bf \bar {z}}}_0 }
\right)\mbox{ of }E\left( {{\rm {\bf z}},{\rm {\bf \bar {z}}}} \right)\mbox{
} \\
 \frac{\partial E\left( {{\rm {\bf z}}_0 ,{\rm {\bf \bar {z}}}_0 }
\right)}{\partial {\rm {\bf x}}}=\frac{1}{S}\left\{ {\frac{\partial H\left(
{{\rm {\bf z}}_0 ,{\rm {\bf \bar {z}}}_0 } \right)}{\partial {\rm {\bf
x}}}-E\frac{\partial S\left( {{\rm {\bf z}}_0 ,{\rm {\bf \bar {z}}}_0 }
\right)}{\partial {\rm {\bf x}}}} \right\}={\rm {\bf 0}}\Rightarrow
\frac{\partial H\left( {{\rm {\bf z}}_0 ,{\rm {\bf \bar {z}}}_0 }
\right)}{\partial {\rm {\bf x}}}=E\frac{\partial S\left( {{\rm {\bf z}}_0
,{\rm {\bf \bar {z}}}_0 } \right)}{\partial {\rm {\bf x}}}. \\
 \mbox{The Hessian of the energy function }E\left( {{\rm {\bf z}},{\rm {\bf
\bar {z}}}} \right)\mbox{ at critical points then becomes } \\
 {\rm {\bf D}}^2E\left( {{\rm {\bf z}}_0 ,{\rm {\bf \bar {z}}}_0 }
\right)=\frac{\partial ^2}{\partial {\rm {\bf x}}\partial {\rm {\bf
y}}}E\left( {{\rm {\bf z}}_0 ,{\rm {\bf \bar {z}}}_0 }
\right)=\frac{1}{S}\left( {\frac{\partial ^2H\left( {{\rm {\bf z}}_0 ,{\rm
{\bf \bar {z}}}_0 } \right)}{\partial {\rm {\bf x}}\partial {\rm {\bf
y}}}-E\left( {{\rm {\bf z}}_0 ,{\rm {\bf \bar {z}}}_0 }
\right)\frac{\partial ^2S\left( {{\rm {\bf z}}_0 ,{\rm {\bf \bar {z}}}_0 }
\right)}{\partial {\rm {\bf x}}\partial {\rm {\bf y}}}} \right)\mbox{; } \\
 \mbox{where }{\rm {\bf x}}=\left( {{\rm {\bf z}},{\rm {\bf \bar {z}}}}
\right)\mbox{ and }{\rm {\bf y}}=\left( {{\rm {\bf z}},{\rm {\bf \bar {z}}}}
\right)\mbox{. } \\
 \end{array}
\]
\underline {\textbf{4. Linear Response.}}

\underline {\textbf{4.1. Exact Linear Response.}}

Exact linear response functions are defined as
\[
\begin{array}{c}
 R_{FA} \left( t \right)=-\frac{i}{\hbar }\theta _+ \left( t
\right)\left\langle {\Psi _0 } \mathrel{\left| {\vphantom {{\Psi _0 }
{\left[ {F\left( t \right),A\left( 0 \right)} \right]\Psi _0 }}} \right.
\kern-\nulldelimiterspace} {\left[ {F\left( t \right),A\left( 0 \right)}
\right]\Psi _0 } \right\rangle \\
 \mbox{ } \\
 \mbox{where }\left| {\Psi _0 } \right\rangle \mbox{ is an exact stationary
state, and the time dependence of observable }A\mbox{ } \\
 \mbox{ is given by }A\left( t \right)=e^{-iH_0 t}Ae^{iH_0 t}\mbox{, where
}H_0 \mbox{ is the unperturbed hamiltonian. } \\
 \end{array}
\]
If $A$ and $F $are one particle operators these functions describe the response of
expectation value of the observable $A$ to the perturbation $F$. This can be
formulated in terms of the first order reduced density matrix elements
response to external time dependent perturbing fields, $F(t) , $that interact with
the system through one body forces, according to the following:
\[
\begin{array}{c}
 F\left( t \right)=\sum\limits_{ij} {f_{kl} \left( t \right)\mbox{ }}
\mbox{a}_k^\dag \mbox{a}_l \\
 \rho _{ij} \left( t \right)=\rho _{ij} \left( 0 \right)+\frac{i}{\hbar
}\int_0^t {\left\langle {\Psi _0 } \mathrel{\left| {\vphantom {{\Psi _0 }
{\left[ {F\left( \tau \right)\mbox{ ,a}_i^\dag \mbox{a}_j } \right]\Psi _0
}}} \right. \kern-\nulldelimiterspace} {\left[ {F\left( \tau \right)\mbox{
,a}_i^\dag \mbox{a}_j } \right]\Psi _0 } \right\rangle } d\tau \\
 \rho _{ij} \left( t \right)=\rho _{ij} \left( 0 \right)+\delta \rho \left(
t \right) \\
 \delta \rho _{ij} \left( t \right)=\sum\limits_{kl} {\int_{-\infty
}^{+\infty } {R_{ijkl} \left( {t-\tau } \right)f_{kl} \left( \tau
\right)d\tau } } \\
 R_{ijkl} \left( t \right)=-\frac{i}{\hbar }\left\langle {\Psi _0 }
\mathrel{\left| {\vphantom {{\Psi _0 } {\left[ {\mbox{a}_j^\dag \left( t
\right)\mbox{a}_i \left( t \right)\mbox{ ,a}_k^\dag \mbox{a}_l } \right]\Psi
_0 }}} \right. \kern-\nulldelimiterspace} {\left[ {\mbox{a}_j^\dag \left( t
\right)\mbox{a}_i \left( t \right)\mbox{ ,a}_k^\dag \mbox{a}_l } \right]\Psi
_0 } \right\rangle \\
 \mbox{where }\rho _{ij} \mbox{ are one particle reduced density matrix
elements and }R_{ijkl} \left( t \right)\mbox{ are } \\
 \mbox{the one particle linear response functions } \\
 \end{array}
\]
The response function may then be written as
\[
\begin{array}{c}
 R_{ijkl} \left( t \right)=-\frac{i}{\hbar }\theta _+ \left( t
\right)\left\{ {\left\langle {\Psi _0 } \mathrel{\left| {\vphantom {{\Psi _0
} {\mbox{a}_j^\dag \mbox{a}_i \mbox{ }e^{-i\left( {H_0 -E_0 }
\right)\raise0.5ex\hbox{$\scriptstyle {t\mbox{
}}$}\kern-0.1em/\kern-0.15em\lower0.25ex\hbox{$\scriptstyle \hbar
$}}\mbox{a}_k^\dag \mbox{a}_l \Psi _0 }}} \right. \kern-\nulldelimiterspace}
{\mbox{a}_j^\dag \mbox{a}_i \mbox{ }e^{-i\left( {H_0 -E_0 }
\right)\raise0.5ex\hbox{$\scriptstyle {t\mbox{
}}$}\kern-0.1em/\kern-0.15em\lower0.25ex\hbox{$\scriptstyle \hbar
$}}\mbox{a}_k^\dag \mbox{a}_l \Psi _0 } \right\rangle -\left\langle {\Psi _0
} \mathrel{\left| {\vphantom {{\Psi _0 } {\mbox{a}_k^\dag \mbox{a}_l \mbox{
}e^{-i\left( {H_0 -E_0 } \right)\raise0.5ex\hbox{$\scriptstyle {t\mbox{
}}$}\kern-0.1em/\kern-0.15em\lower0.25ex\hbox{$\scriptstyle \hbar
$}}\mbox{a}_j^\dag \mbox{a}_i \Psi _0 }}} \right. \kern-\nulldelimiterspace}
{\mbox{a}_k^\dag \mbox{a}_l \mbox{ }e^{-i\left( {H_0 -E_0 }
\right)\raise0.5ex\hbox{$\scriptstyle {t\mbox{
}}$}\kern-0.1em/\kern-0.15em\lower0.25ex\hbox{$\scriptstyle \hbar
$}}\mbox{a}_j^\dag \mbox{a}_i \Psi _0 } \right\rangle } \right\} \\
 \mbox{where }E_0 \mbox{ is the energy of the stationary state }\left| {\Psi
_0 } \right\rangle \mbox{. Taking the Laplace transform } \\
 \mbox{of }R_{ijkl} \left( t \right)\mbox{ we obtain } \\
 R_{ijkl} \left( z \right)=\sum\limits_{n\ne 0} {\left\{ {\frac{\left\langle
{\Psi _0 } \mathrel{\left| {\vphantom {{\Psi _0 } {\mbox{a}_j^\dag
\mbox{a}_i \Psi _n }}} \right. \kern-\nulldelimiterspace} {\mbox{a}_j^\dag
\mbox{a}_i \Psi _n } \right\rangle \left\langle {\Psi _n } \mathrel{\left|
{\vphantom {{\Psi _n } {\mbox{a}_k^\dag \mbox{a}_l \Psi _0 }}} \right.
\kern-\nulldelimiterspace} {\mbox{a}_k^\dag \mbox{a}_l \Psi _0 }
\right\rangle }{z-\left( {E_n -E_0 } \right)}-\frac{\left\langle {\Psi _0 }
\mathrel{\left| {\vphantom {{\Psi _0 } {\mbox{a}_k^\dag \mbox{a}_l \Psi _n
}}} \right. \kern-\nulldelimiterspace} {\mbox{a}_k^\dag \mbox{a}_l \Psi _n }
\right\rangle \left\langle {\Psi _n } \mathrel{\left| {\vphantom {{\Psi _n }
{\mbox{a}_j^\dag \mbox{a}_i \Psi _0 }}} \right. \kern-\nulldelimiterspace}
{\mbox{a}_j^\dag \mbox{a}_i \Psi _0 } \right\rangle }{z+\left( {E_n -E_0 }
\right)}} \right\}} \\
 \mbox{The response functions }R_{ijkl} \left( z \right)\mbox{ form a
hermitian matrix for real }z,\mbox{ thus } \\
 \mbox{can be diagonalized to give } \\
 \tilde {R}_{KL} \left( z \right)=g_K \left( z \right)\delta _{KL}
=\sum\limits_{n\ne 0} {\left\{ {\frac{\left\langle {\Psi _0 }
\mathrel{\left| {\vphantom {{\Psi _0 } {T_K \Psi _n }}} \right.
\kern-\nulldelimiterspace} {T_K \Psi _n } \right\rangle \left\langle {\Psi
_n } \mathrel{\left| {\vphantom {{\Psi _n } {T_K^\dag \Psi _0 }}} \right.
\kern-\nulldelimiterspace} {T_K^\dag \Psi _0 } \right\rangle }{z-\left( {E_n
-E_0 } \right)}-\frac{\left\langle {\Psi _0 } \mathrel{\left| {\vphantom
{{\Psi _0 } {T_K \Psi _n }}} \right. \kern-\nulldelimiterspace} {T_K \Psi _n
} \right\rangle \left\langle {\Psi _n } \mathrel{\left| {\vphantom {{\Psi _n
} {T_K^\dag \Psi _0 }}} \right. \kern-\nulldelimiterspace} {T_K^\dag \Psi _0
} \right\rangle }{z+\left( {E_n -E_0 } \right)}} \right\}} \\
 \mbox{where }T_K^\dag =\sum\limits_{ij} {\mbox{T}_{Kij} \mbox{a}_i^\dag
\mbox{a}_j } \mbox{ } \\
 \end{array}
\]
The poles of the Laplace transform occur at $\pm \left( {E_n -E_0 } \right)$
and the residues $\left| {\left\langle {\Psi _n } \mathrel{\left| {\vphantom
{{\Psi _n } {T_K^\dag \Psi _0 }}} \right. \kern-\nulldelimiterspace}
{T_K^\dag \Psi _0 } \right\rangle } \right|^2$ are transition probabilities
that the operator $T_K^\dag $ will cause a transition from the stationary
state $\left| {\Psi _0 } \right\rangle $ to the stationary state $\left|
{\Psi _n } \right\rangle $.

In general the excited states are of the form
\[
\left| {\Psi _n } \right\rangle =\sum\limits_{m=0}^{m=N} {\sum\limits_{K_1
K_2 \cdots K_m } {c_{K_1 K_2 \cdots K_m }^n _ T_{K_1 }^\dag T_{K_2 }^\dag
\cdots T_{K_m }^\dag } } \left| {\Psi _0 } \right\rangle
\]
In this study we focus on excited states that are produced from the
reference state by single photon processes, i.e. low intensity
electromagnetic radiation, which is consistent with the hypotheses of linear
response. Thus we desire that the excited states be well approximated by
\[
\left| {\Psi _n } \right\rangle \approx T_n^\dag \left| {\Psi _0 }
\right\rangle
\]
\underline {\textbf{4.2. Linear Response in a Coherent State Setting}}

The classical Hamiltonian corresponding to the Hamiltonian containing the
time dependent perturbation is
\[
E\left( {{\rm {\bf z}},{\rm {\bf \bar {z}}},t} \right)=\frac{\left\langle
{{\rm {\bf z}}} \mathrel{\left| {\vphantom {{\rm {\bf z}} {\left( {H_0
+F\left( t \right)} \right){\rm {\bf z}}}}} \right.
\kern-\nulldelimiterspace} {\left( {H_0 +F\left( t \right)} \right){\rm {\bf
z}}} \right\rangle }{\left\langle {{\rm {\bf z}}} \mathrel{\left| {\vphantom
{{\rm {\bf z}} {\rm {\bf z}}}} \right. \kern-\nulldelimiterspace} {{\rm {\bf
z}}} \right\rangle }=E_0 \left( {{\rm {\bf z}},{\rm {\bf \bar {z}}}}
\right)+\frac{\left\langle {{\rm {\bf z}}} \mathrel{\left| {\vphantom {{\rm
{\bf z}} {F\left( t \right){\rm {\bf z}}}}} \right.
\kern-\nulldelimiterspace} {F\left( t \right){\rm {\bf z}}} \right\rangle
}{\left\langle {{\rm {\bf z}}} \mathrel{\left| {\vphantom {{\rm {\bf z}}
{\rm {\bf z}}}} \right. \kern-\nulldelimiterspace} {{\rm {\bf z}}}
\right\rangle }
\]
One can form a quadratic approximation to this function around a point
$\left( {{\rm {\bf z}}_0 ,{\rm {\bf \bar {z}}}_0 } \right)\mbox{ }$ by
expanding up to second order
\[
\begin{array}{c}
 E\left( {{\rm {\bf z}}_{\rm {\bf 0}} +\delta {\rm {\bf z}},{\rm {\bf \bar
{z}}}_{\rm {\bf 0}} +\delta {\rm {\bf \bar {z}}},t} \right)\approx E_2
\left( {{\rm {\bf z}}_{\rm {\bf 0}} ,{\rm {\bf \bar {z}}}_{\rm {\bf 0}} ,t}
\right)=E\left( {{\rm {\bf z}}_{\rm {\bf 0}} ,{\rm {\bf \bar {z}}}_{\rm {\bf
0}} ,t} \right)+\delta ^1E\left( {{\rm {\bf z}}_{\rm {\bf 0}} ,{\rm {\bf
\bar {z}}}_{\rm {\bf 0}} ,t} \right)+\delta ^2E\left( {{\rm {\bf z}}_{\rm
{\bf 0}} ,{\rm {\bf \bar {z}}}_{\rm {\bf 0}} ,t} \right) \\
 \mbox{where }\delta ^1E\left( {{\rm {\bf z}}_{\rm {\bf 0}} ,{\rm {\bf \bar
{z}}}_{\rm {\bf 0}} ,t} \right)=\left( {\delta {\rm {\bf z}},\delta {\rm
{\bf \bar {z}}}} \right){\rm {\bf D}}^1E\left( {{\rm {\bf z}}_0 ,{\rm {\bf
\bar {z}}}_0 ,t} \right)\mbox{ and }\delta ^2E\left( {{\rm {\bf z}}_{\rm
{\bf 0}} ,{\rm {\bf \bar {z}}}_{\rm {\bf 0}} ,t} \right)=\frac{1}{2}\left(
{\delta {\rm {\bf \bar {z}}},\delta {\rm {\bf z}}} \right){\rm {\bf
D}}^2E\left( {{\rm {\bf z}}_0 ,{\rm {\bf \bar {z}}}_0 ,t} \right)\left(
{{\begin{array}{*{20}c}
 {\delta {\rm {\bf \bar {z}}}} \hfill \\
 {\delta {\rm {\bf z}}} \hfill \\
\end{array} }} \right)\mbox{ } \\
 \end{array}
\]
The approximate classical Hamiltonian associated with linear response about
the point $\left( {{\rm {\bf z}}_0 ,{\rm {\bf \bar {z}}}_0 } \right)$ is
then constructed by only including the time dependent perturbation in terms
up to first order in $\left( {\delta {\rm {\bf z}},\delta {\rm {\bf \bar
{z}}}} \right)$\underline { }
\[
\tilde {E}\left( {{\rm {\bf z}}_0 +\delta {\rm {\bf z}},{\rm {\bf \bar
{z}}}_0 +\delta {\rm {\bf \bar {z}}}} \right)=E\left( {{\rm {\bf z}}_0 ,{\rm
{\bf \bar {z}}}_0 } \right)+\left( {{\begin{array}{*{20}c}
 {\frac{\partial E\left( {{\rm {\bf z}}_0 ,{\rm {\bf \bar {z}}}_0 ,t}
\right)}{\partial {\rm {\bf \bar {z}}}}} \hfill & {\frac{\partial E\left(
{{\rm {\bf z}}_0 ,{\rm {\bf \bar {z}}}_0 ,t} \right)}{\partial {\rm {\bf
z}}}} \hfill \\
\end{array} }} \right)\left( {{\begin{array}{*{20}c}
 {\delta {\rm {\bf \bar {z}}}} \hfill \\
 {\delta {\rm {\bf z}}} \hfill \\
\end{array} }} \right)+\frac{1}{2}\left( {{\begin{array}{*{20}c}
 {\frac{\partial ^2E_0 \left( {{\rm {\bf z}}_0 ,{\rm {\bf \bar {z}}}_0 }
\right)}{\partial {\rm {\bf \bar {z}}}\partial {\rm {\bf \bar {z}}}}} \hfill
& {\frac{\partial ^2E_0 \left( {{\rm {\bf z}}_0 ,{\rm {\bf \bar {z}}}_0 }
\right)}{\partial {\rm {\bf \bar {z}}}\partial {\rm {\bf z}}}} \hfill \\
 {\frac{\partial ^2E_0 \left( {{\rm {\bf z}}_0 ,{\rm {\bf \bar {z}}}_0 }
\right)}{\partial {\rm {\bf z}}\partial {\rm {\bf \bar {z}}}}} \hfill &
{\frac{\partial ^2E_0 \left( {{\rm {\bf z}}_0 ,{\rm {\bf \bar {z}}}_0 }
\right)}{\partial {\rm {\bf z}}\partial {\rm {\bf z}}}} \hfill \\
\end{array} }} \right)
\]
\[
\begin{array}{c}
 \mbox{Using the linear response approximation one obtains an approximate
gradient for } \\
 \mbox{local variations } \\
 \left( {{\begin{array}{*{20}c}
 {\frac{\partial \tilde {E}_2 }{\partial \delta {\rm {\bf \bar {z}}}}}
\hfill \\
 {\frac{\partial \tilde {E}_2 }{\partial \delta {\rm {\bf z}}}} \hfill \\
\end{array} }} \right)={\rm {\bf D}}^1E\left( {{\rm {\bf \bar {z}}}_0 ,{\rm
{\bf z}}_0 ,t} \right)+{\rm {\bf D}}^2E_0 \left( {{\rm {\bf \bar {z}}}_0
,{\rm {\bf z}}_0 } \right)\left( {{\begin{array}{*{20}c}
 {\delta {\rm {\bf \bar {z}}}} \hfill \\
 {\delta {\rm {\bf z}}} \hfill \\
\end{array} }} \right) \\
 \mbox{which can be used to obtain a linearized EOM in the tangent space at
the point }\left( {{\rm {\bf z}}_0 ,{\rm {\bf \bar {z}}}_0 } \right) \\
 \left( {{\begin{array}{*{20}c}
 {\rm {\bf 0}} \hfill & {\tfrac{i}{\hbar }{\rm {\bf C}}} \hfill \\
 {-\tfrac{i}{\hbar }{\rm {\bf \bar {C}}}} \hfill & {\rm {\bf 0}} \hfill \\
\end{array} }} \right)\left( {{\begin{array}{*{20}c}
 {\delta {\rm {\bf \dot {\bar {z}}}}} \hfill \\
 {\delta {\rm {\bf \dot {z}}}} \hfill \\
\end{array} }} \right)={\rm {\bf D}}^1E\left( {{\rm {\bf \bar {z}}}_0 ,{\rm
{\bf z}}_0 ,t} \right)+{\rm {\bf D}}^2E_0 \left( {{\rm {\bf \bar {z}}}_0
,{\rm {\bf z}}_0 } \right)\left( {{\begin{array}{*{20}c}
 {\delta {\rm {\bf \bar {z}}}} \hfill \\
 {\delta {\rm {\bf z}}} \hfill \\
\end{array} }} \right) \\
 \mbox{ } \\
 \end{array}
\]
If the point $\left( {{\rm {\bf z}}_0 ,{\rm {\bf \bar {z}}}_0 } \right)$ is
a critical point of the unperturbed energy $E_0 \left( {{\rm {\bf z}},{\rm
{\bf \bar {z}}}} \right)$ then
\[
\begin{array}{c}
 \frac{\partial E\left( {{\rm {\bf z}}_0 ,{\rm {\bf \bar {z}}}_0 ,t}
\right)}{\partial {\rm {\bf x}}}=\frac{1}{S}\left\{ {\frac{\partial F\left(
{{\rm {\bf z}}_0 ,{\rm {\bf \bar {z}}}_0 ,t} \right)}{\partial {\rm {\bf
x}}}-\frac{\left\langle {{\rm {\bf z}}_{\rm {\bf 0}} } \mathrel{\left|
{\vphantom {{{\rm {\bf z}}_{\rm {\bf 0}} } {F\left( 0 \right){\rm {\bf
z}}_{\rm {\bf 0}} }}} \right. \kern-\nulldelimiterspace} {F\left( 0
\right){\rm {\bf z}}_{\rm {\bf 0}} } \right\rangle }{\left\langle {{\rm {\bf
z}}_{\rm {\bf 0}} } \mathrel{\left| {\vphantom {{{\rm {\bf z}}_{\rm {\bf 0}}
} {{\rm {\bf z}}_{\rm {\bf 0}} }}} \right. \kern-\nulldelimiterspace} {{\rm
{\bf z}}_{\rm {\bf 0}} } \right\rangle }\frac{\partial S\left( {{\rm {\bf
z}}_0 ,{\rm {\bf \bar {z}}}_0 ,t} \right)}{\partial {\rm {\bf x}}}}
\right\}=\frac{1}{S}\frac{\partial F\left( {{\rm {\bf z}}_0 ,{\rm {\bf \bar
{z}}}_0 ,t} \right)}{\partial {\rm {\bf x}}} \\
 \mbox{as }{\rm {\bf z}}_0 ={\rm {\bf z}}\left( 0 \right)\mbox{ and at
}t\mbox{=}0\mbox{ the external perturbation }F=F\left( 0 \right)=0\mbox{ .
The Hessian at this point } \\
 \frac{\partial ^2E_0 \left( {{\rm {\bf z}}_0 ,{\rm {\bf \bar {z}}}_0 ,t}
\right)}{\partial {\rm {\bf x}}\partial {\rm {\bf y}}}=\frac{1}{S}\left\{
{\frac{\partial ^2H_0 \left( {{\rm {\bf z}}_0 ,{\rm {\bf \bar {z}}}_0 ,t}
\right)}{\partial {\rm {\bf x}}\partial {\rm {\bf y}}}-E\left( {{\rm {\bf
z}}_0 ,{\rm {\bf \bar {z}}}_0 } \right)\frac{\partial ^2S\left( {{\rm {\bf
z}}_0 ,{\rm {\bf \bar {z}}}_0 ,t} \right)}{\partial {\rm {\bf x}}\partial
{\rm {\bf y}}}} \right\} \\
 \mbox{where }{\rm {\bf x}},{\rm {\bf y}}={\rm {\bf z}},{\rm {\bf \bar
{z}}}\mbox{ } \\
 \end{array}
\]
The linear response $\left( {\delta {\rm {\bf \bar {z}}},\delta {\rm {\bf
z}}} \right)$  of the state around a critical point is then determined by the
equation
\[
\begin{array}{c}
 \left( {{\begin{array}{*{20}c}
 {\rm {\bf 0}} \hfill & {\tfrac{i}{\hbar }{\rm {\bf C}}} \hfill \\
 {-\tfrac{i}{\hbar }{\rm {\bf \bar {C}}}} \hfill & {\rm {\bf 0}} \hfill \\
\end{array} }} \right)\left( {{\begin{array}{*{20}c}
 {\delta {\rm {\bf \dot {\bar {z}}}}} \hfill \\
 {\delta {\rm {\bf \dot {z}}}} \hfill \\
\end{array} }} \right)=\left( {{\begin{array}{*{20}c}
 {\frac{\partial \tilde {E}\left( {{\rm {\bf z}}_0 ,{\rm {\bf \bar {z}}}_0 }
\right)}{\partial \delta {\rm {\bf \bar {z}}}}} \hfill \\
 {\frac{\partial \tilde {E}\left( {{\rm {\bf z}}_0 ,{\rm {\bf \bar {z}}}_0 }
\right)}{\partial \delta {\rm {\bf z}}}} \hfill \\
\end{array} }} \right)={\rm {\bf D}}^{\rm {\bf 1}}F\left( {{\rm {\bf z}}_0
,{\rm {\bf \bar {z}}}_0 ,t} \right)+{\rm {\bf D}}^2E_0 \left( {{\rm {\bf
z}}_0 ,{\rm {\bf \bar {z}}}_0 } \right)\left( {{\begin{array}{*{20}c}
 {\delta {\rm {\bf \bar {z}}}} \hfill \\
 {\delta {\rm {\bf z}}} \hfill \\
\end{array} }} \right) \\
 \Rightarrow \left( {\left( {{\begin{array}{*{20}c}
 {\rm {\bf 0}} \hfill & {\tfrac{i}{\hbar }{\rm {\bf C}}\frac{{\rm {\bf
d}}}{{\rm {\bf dt}}}} \hfill \\
 {-\tfrac{i}{\hbar }{\rm {\bf \bar {C}}}\frac{{\rm {\bf d}}}{{\rm {\bf
dt}}}} \hfill & {\rm {\bf 0}} \hfill \\
\end{array} }} \right)-{\rm {\bf D}}^2E_0 \left( {{\rm {\bf z}}_0 ,{\rm {\bf
\bar {z}}}_0 } \right)} \right)\left( {{\begin{array}{*{20}c}
 {\delta {\rm {\bf \bar {z}}}} \hfill \\
 {\delta {\rm {\bf z}}} \hfill \\
\end{array} }} \right)={\rm {\bf D}}^{\rm {\bf 1}}F\left( {{\rm {\bf z}}_0
,{\rm {\bf \bar {z}}}_0 ,t} \right) \\
 \Rightarrow \left( {{\begin{array}{*{20}c}
 {\delta {\rm {\bf \bar {z}}}} \hfill \\
 {\delta {\rm {\bf z}}} \hfill \\
\end{array} }} \right)\left( t \right)=\int\limits_{-\infty }^\infty
{R\left( {t-\tau } \right){\rm {\bf D}}^{\rm {\bf 1}}F\left( {{\rm {\bf
z}}_0 ,{\rm {\bf \bar {z}}}_0 ,\tau } \right)} d\tau \\
 \mbox{where the response function is defined as } \\
 R\left( t \right)=\theta _+ \left( t \right)\mbox{ }\left( {\left(
{{\begin{array}{*{20}c}
 {\rm {\bf 0}} \hfill & {\tfrac{i}{\hbar }{\rm {\bf C}}\frac{{\rm {\bf
d}}}{{\rm {\bf dt}}}} \hfill \\
 {-\tfrac{i}{\hbar }{\rm {\bf \bar {C}}}\frac{{\rm {\bf d}}}{{\rm {\bf
dt}}}} \hfill & {\rm {\bf 0}} \hfill \\
\end{array} }} \right)-{\rm {\bf D}}^2E_0 \left( {{\rm {\bf z}}_0 ,{\rm {\bf
\bar {z}}}_0 } \right)} \right)^{-1}\left( t \right)\mbox{ } \\
 \mbox{Taking the Laplace transform we obtain } \\
 \left( {{\begin{array}{*{20}c}
 {\delta {\rm {\bf \bar {z}}}} \hfill \\
 {\delta {\rm {\bf z}}} \hfill \\
\end{array} }} \right)\left( z \right)=R\left( z \right){\rm {\bf D}}^{\rm
{\bf 1}}F\left( {{\rm {\bf z}}_0 ,{\rm {\bf \bar {z}}}_0 ,z} \right)\mbox{,
where }R\left( z \right)\mbox{= }\left( {\left( {{\begin{array}{*{20}c}
 {\rm {\bf 0}} \hfill & {-z{\rm {\bf C}}} \hfill \\
 {z{\rm {\bf \bar {C}}}} \hfill & {\rm {\bf 0}} \hfill \\
\end{array} }} \right)-{\rm {\bf D}}^2E_0 \left( {{\rm {\bf z}}_0 ,{\rm {\bf
\bar {z}}}_0 } \right)} \right)^{-1}\mbox{ } \\
 \end{array}
\]
The preceding response function $R\left( t \right)$ describes the response
of observables of a system described by $H_0 $ to external perturbations
$F\left( t \right)$ that interact through 1,2,....,N body forces. The
response of the observables are described by the response of the density and
reduced density matrices $\left\{ {\rho ^\mbox{N}\left( {{\rm {\bf z}}_0
,{\rm {\bf z}}_0 } \right),\cdots ,\rho ^\mbox{1}\left( {{\rm {\bf z}}_0
,{\rm {\bf z}}_0 } \right)} \right\}$. The linear response of the matrix
elements $\rho ^\mbox{N}\left( {{\rm {\bf z}}_0 ,{\rm {\bf \bar {z}}}_0 }
\right)_{KL} $ of the N-particle density operator $\left| {{\rm {\bf z}}_0 }
\right\rangle \left\langle {{\rm {\bf z}}_0 } \right|$ of the coherent state
$\left| {{\rm {\bf z}}_0 } \right\rangle $ to the external perturbation are
given by
\[
\delta ^1\left\{ {\left\langle {\Phi _K } \mathrel{\left| {\vphantom {{\Phi
_K } {{\rm {\bf z}}_0 }}} \right. \kern-\nulldelimiterspace} {{\rm {\bf
z}}_0 } \right\rangle \left\langle {{\rm {\bf z}}_0 } \mathrel{\left|
{\vphantom {{{\rm {\bf z}}_0 } {\Phi _L }}} \right.
\kern-\nulldelimiterspace} {\Phi _L } \right\rangle } \right\}\left( t
\right)=\delta ^1\left\{ {\left\langle {{\rm {\bf z}}_0 } \mathrel{\left|
{\vphantom {{{\rm {\bf z}}_0 } {\Phi _L }}} \right.
\kern-\nulldelimiterspace} {\Phi _L } \right\rangle \left\langle {\Phi _K }
\mathrel{\left| {\vphantom {{\Phi _K } {{\rm {\bf z}}_0 }}} \right.
\kern-\nulldelimiterspace} {{\rm {\bf z}}_0 } \right\rangle }
\right\}=\delta ^1\Phi _{KL} \left( {{\rm {\bf z}}_0 ,{\rm {\bf \bar {z}}}_0
} \right)=\left( {{\begin{array}{*{20}c}
 {\frac{\partial \Phi _{KL} \left( {{\rm {\bf z}}_0 ,{\rm {\bf \bar {z}}}_0
} \right)}{\partial {\rm {\bf \bar {z}}}}} \hfill & {\frac{\partial \Phi
_{KL} \left( {{\rm {\bf z}}_0 ,{\rm {\bf \bar {z}}}_0 } \right)}{\partial
{\rm {\bf z}}}} \hfill \\
\end{array} }} \right)\left( {{\begin{array}{*{20}c}
 {\delta {\rm {\bf \bar {z}}}} \hfill \\
 {\delta {\rm {\bf z}}} \hfill \\
\end{array} }} \right)
\]
which have Laplace transforms
\[
\left( {{\begin{array}{*{20}c}
 {\frac{\partial \Phi _{KL} \left( {{\rm {\bf z}}_0 ,{\rm {\bf \bar {z}}}_0
} \right)}{\partial {\rm {\bf \bar {z}}}}} \hfill & {\frac{\partial \Phi
_{KL} \left( {{\rm {\bf z}}_0 ,{\rm {\bf \bar {z}}}_0 } \right)}{\partial
{\rm {\bf z}}}} \hfill \\
\end{array} }} \right)\left( {\left( {{\begin{array}{*{20}c}
 {\rm {\bf 0}} \hfill & {-z{\rm {\bf C}}} \hfill \\
 {z{\rm {\bf \bar {C}}}} \hfill & {\rm {\bf 0}} \hfill \\
\end{array} }} \right)-{\rm {\bf D}}^2E_0 \left( {{\rm {\bf z}}_0 ,{\rm {\bf
\bar {z}}}_0 } \right)} \right)^{-1}{\rm {\bf D}}^{\rm {\bf 1}}F\left( {{\rm
{\bf z}}_0 ,{\rm {\bf \bar {z}}}_0 ,z} \right).
\]
The matrix elements and the linear response of the 1-particle reduced
density operator can be constructed in the following manner:
\[
\begin{array}{c}
 \rho ^\mbox{N}\left( {{\rm {\bf z}}_0 ,{\rm {\bf \bar {z}}}_0 }
\right)=\sum\limits_{K,L} {\Phi _{KL} \left( {{\rm {\bf z}}_0 ,{\rm {\bf
\bar {z}}}_0 } \right)\left| {\Phi _L } \right\rangle \left\langle {\Phi _K
} \right|} \Rightarrow \rho ^1\left( {{\rm {\bf z}}_0 ,{\rm {\bf \bar
{z}}}_0 } \right)_{ij} =Tr\left\{ {\rho ^\mbox{N}\left( {{\rm {\bf z}}_0
,{\rm {\bf \bar {z}}}_0 } \right)\mbox{a}_j^\dag \mbox{a}_i } \right\} \\
 \therefore \delta ^1\rho ^\mbox{N}\left( {{\rm {\bf z}}_0 ,{\rm {\bf \bar
{z}}}_0 } \right)=\sum\limits_{K,L} {\delta ^1\Phi _{KL} \left( {{\rm {\bf
z}}_0 ,{\rm {\bf \bar {z}}}_0 } \right)\left| {\Phi _L } \right\rangle
\left\langle {\Phi _K } \right|} \Rightarrow \delta ^1\rho ^1\left( {{\rm
{\bf z}}_0 ,{\rm {\bf \bar {z}}}_0 } \right)_{ij} =Tr\left\{ {\delta ^1\rho
^\mbox{N}\left( {{\rm {\bf z}}_0 ,{\rm {\bf \bar {z}}}_0 }
\right)\mbox{a}_j^\dag \mbox{a}_i } \right\} \\
 \Rightarrow \delta ^1\rho ^\mbox{1}\left( {{\rm {\bf z}}_0 ,{\rm {\bf \bar
{z}}}_0 } \right)=Tr\left\{ {\delta ^1\rho ^\mbox{N}\left( {{\rm {\bf z}}_0
,{\rm {\bf \bar {z}}}_0 } \right)\mbox{a}_j^\dag \mbox{a}_i }
\right\}=Tr\left\{ {\sum\limits_{K,L} {\delta ^1\Phi _{KL} \left( {{\rm {\bf
z}}_0 ,{\rm {\bf \bar {z}}}_0 } \right)\left| {\Phi _L } \right\rangle
\left\langle {\Phi _K } \right|} \mbox{a}_j^\dag \mbox{a}_i } \right\} \\
 =Tr\left\{ {\sum\limits_{K,L} {\left( {{\begin{array}{*{20}c}
 {\frac{\partial \Phi _{KL} \left( {{\rm {\bf z}}_0 ,{\rm {\bf \bar {z}}}_0
} \right)}{\partial {\rm {\bf \bar {z}}}}} \hfill & {\frac{\partial \Phi
_{KL} \left( {{\rm {\bf z}}_0 ,{\rm {\bf \bar {z}}}_0 } \right)}{\partial
{\rm {\bf z}}}} \hfill \\
\end{array} }} \right)\left( {{\begin{array}{*{20}c}
 {\delta {\rm {\bf \bar {z}}}} \hfill \\
 {\delta {\rm {\bf z}}} \hfill \\
\end{array} }} \right)\left| {\Phi _L } \right\rangle \left\langle {\Phi _K
} \right|} \mbox{a}_j^\dag \mbox{a}_i } \right\} \\
 =\left( {{\begin{array}{*{20}c}
 {\frac{\partial Tr\left\{ {\sum\limits_{K,L} {\Phi _{KL} \left( {{\rm {\bf
z}}_0 ,{\rm {\bf \bar {z}}}_0 } \right)\left| {\Phi _L } \right\rangle
\left\langle {\Phi _K } \right|\mbox{a}_j^\dag \mbox{a}_i } }
\right\}}{\partial {\rm {\bf \bar {z}}}}} \hfill & {\frac{\partial Tr\left\{
{\sum\limits_{K,L} {\Phi _{KL} \left( {{\rm {\bf z}}_0 ,{\rm {\bf \bar
{z}}}_0 } \right)\left| {\Phi _L } \right\rangle \left\langle {\Phi _K }
\right|\mbox{a}_j^\dag \mbox{a}_i } } \right\}}{\partial {\rm {\bf z}}}}
\hfill \\
\end{array} }} \right)\left( {{\begin{array}{*{20}c}
 {\delta {\rm {\bf \bar {z}}}} \hfill \\
 {\delta {\rm {\bf z}}} \hfill \\
\end{array} }} \right) \\
 =\left( {{\begin{array}{*{20}c}
 {\frac{\partial \rho ^1\left( {{\rm {\bf z}}_0 ,{\rm {\bf \bar {z}}}_0 }
\right)_{ij} }{\partial {\rm {\bf \bar {z}}}}} \hfill & {\frac{\partial \rho
^1\left( {{\rm {\bf z}}_0 ,{\rm {\bf \bar {z}}}_0 } \right)_{ij} }{\partial
{\rm {\bf z}}}} \hfill \\
\end{array} }} \right)\left( {{\begin{array}{*{20}c}
 {\delta {\rm {\bf \bar {z}}}} \hfill \\
 {\delta {\rm {\bf z}}} \hfill \\
\end{array} }} \right) \\
 \end{array}
\]
The Laplace transform of this response is then
\[
\delta ^1\rho ^1\left( {{\rm {\bf z}}_0 ,{\rm {\bf \bar {z}}}_0 }
\right)_{ij} \left( z \right)=\left( {{\begin{array}{*{20}c}
 {\frac{\partial \rho ^1\left( {{\rm {\bf z}}_0 ,{\rm {\bf \bar {z}}}_0 }
\right)_{ij} }{\partial {\rm {\bf \bar {z}}}}} \hfill & {\frac{\partial \rho
^1\left( {{\rm {\bf z}}_0 ,{\rm {\bf \bar {z}}}_0 } \right)_{ij} }{\partial
{\rm {\bf z}}}} \hfill \\
\end{array} }} \right)\left( {\left( {{\begin{array}{*{20}c}
 {\rm {\bf 0}} \hfill & {-z{\rm {\bf C}}} \hfill \\
 {z{\rm {\bf \bar {C}}}} \hfill & {\rm {\bf 0}} \hfill \\
\end{array} }} \right)-{\rm {\bf D}}^2E_0 \left( {{\rm {\bf z}}_0 ,{\rm {\bf
\bar {z}}}_0 } \right)} \right)^{-1}{\rm {\bf D}}^{\rm {\bf 1}}F\left( {{\rm
{\bf z}}_0 ,{\rm {\bf \bar {z}}}_0 ,z} \right)
\]
\underline {\textbf{4.3 Exact Linear Response in Coherent State Setting}}

In order to facilitate a direct comparison of approximations to exact linear
response functions it is useful to derive the exact linear response function
in terms of coherent states based on the coset
$\raise0.7ex\hbox{${\mbox{SU}\left( {H^\mbox{N}\mbox{ }} \right)}$}
\!\mathord{\left/ {\vphantom {{\mbox{SU}\left( {H^\mbox{N}\mbox{ }} \right)}
{\mbox{U}\left( \mbox{1}
\right)}}}\right.\kern-\nulldelimiterspace}\!\lower0.7ex\hbox{${\mbox{U}\left(
\mbox{1} \right)}$}$. The classical equations of motion defined on this
coherent state manifold are equivalent to the full time dependent
Schroedinger equation. The configuration space formed by these coherent
sates is constituted by all of the states $\left\{ {\left| {\rm {\bf z}}
\right\rangle } \right\}\mbox{ in }H^\mbox{N} $. The group
$\mbox{G=SU}\left( {H^\mbox{N}} \right)$ and the subgroup H is defined by
\[
\begin{array}{c}
 \mbox{H=}\left\{ {\mbox{g}\in \mbox{G}\mathrel\backepsilon \mbox{g}\left|
{\Phi _1 } \right\rangle =\left| {\Phi _1 } \right\rangle }
\right\}=e^{\begin{array}{l}
 i\lambda \left| {\Phi _1 } \right\rangle \left\langle {\Phi _1 } \right| \\
 \\
 \end{array}}e^{i\sum\limits_{ij\ne 1} {\lambda _{ij} \left| {\Phi _i }
\right\rangle \left\langle {\Phi _j } \right|} } \\
 \mbox{to produce coset representatives } \\
 \mbox{G/H=}e^{\sum\limits_{i>1} {\eta _i \left| {\Phi _i } \right\rangle
\left\langle {\Phi _1 } \right|} }\mbox{ }e^{\begin{array}{l}
 \eta _1 \left| {\Phi _1 } \right\rangle \left\langle {\Phi _1 } \right| \\
 \\
 \end{array}}\mbox{ }e^{-\sum\limits_{i>1} {\bar {\eta }_i \left| {\Phi _1 }
\right\rangle \left\langle {\Phi _i } \right|} } \\
 \mbox{where }\left\{ {\Phi _i ;1\leqslant i\leqslant \left(
{{\begin{array}{*{20}c}
 r \hfill \\
 N \hfill \\
\end{array} }} \right)} \right\}\mbox{is a complete orthonormal basis for
}H^\mbox{N}.\mbox{ } \\
 \mbox{The family of coherent states associated with this coset is defined
as } \\
 \left| \eta \right\rangle \mbox{ =}e^{\sum\limits_{i>1} {\eta _i \left|
{\Phi _i } \right\rangle \left\langle {\Phi _1 } \right|} }\mbox{
}e^{\begin{array}{l}
 \eta _1 \left| {\Phi _1 } \right\rangle \left\langle {\Phi _1 } \right| \\
 \\
 \end{array}}\left| {\Phi _1 } \right\rangle \\
 \mbox{For convenience we introduce a new cooridnate system}\left\{ {\rm
{\bf z}} \right\}\mbox{ by setting } \\
 \mbox{ }\eta _1 \mbox{=ln}\left( {z_1 } \right)\mbox{ and }\left\{ {\eta _i
=\frac{z_i }{z_1 };\mbox{ }2\leqslant i\leqslant \left(
{{\begin{array}{*{20}c}
 r \hfill \\
 N \hfill \\
\end{array} }} \right)} \right\},\mbox{ } \\
 \mbox{which leads to the following form for the coherent states } \\
 \left| {\rm {\bf z}} \right\rangle =\sum\limits_{1\leqslant K\leqslant
\left( {{\begin{array}{*{20}c}
 r \hfill \\
 N \hfill \\
\end{array} }} \right)} {z_K \left| {\Phi _K } \right\rangle } \mbox{ } \\
 \end{array}
\]
The overlap kernel, classical Hamiltonian and Kaehler metric for this coset
manifold are then defined for an arbitrary quantum Hamiltonian $H$ as:
\[
\begin{array}{c}
 S=S\left( {{\rm {\bf z}},{\rm {\bf z}}} \right)=\left\langle {{\rm {\bf
z}}} \mathrel{\left| {\vphantom {{\rm {\bf z}} {\rm {\bf z}}}} \right.
\kern-\nulldelimiterspace} {{\rm {\bf z}}} \right\rangle
=\sum\limits_{1\leqslant i\leqslant \left( {{\begin{array}{*{20}c}
 r \hfill \\
 N \hfill \\
\end{array} }} \right)} {\left| {z_i } \right|} ^2 \\
 E\left( {{\rm {\bf z}},{\rm {\bf z}}} \right)=\frac{\left\langle {{\rm {\bf
z}}} \mathrel{\left| {\vphantom {{\rm {\bf z}} {H{\rm {\bf z}}}}} \right.
\kern-\nulldelimiterspace} {H{\rm {\bf z}}} \right\rangle }{\left\langle
{{\rm {\bf z}}} \mathrel{\left| {\vphantom {{\rm {\bf z}} {\rm {\bf z}}}}
\right. \kern-\nulldelimiterspace} {{\rm {\bf z}}} \right\rangle
}=\frac{\sum\limits_{1\leqslant i,j\leqslant \left( {{\begin{array}{*{20}c}
 r \hfill \\
 N \hfill \\
\end{array} }} \right)} {\bar {z}_i z_j \left\langle {\Phi _i }
\mathrel{\left| {\vphantom {{\Phi _i } {H\Phi _j }}} \right.
\kern-\nulldelimiterspace} {H\Phi _j } \right\rangle }
}{\sum\limits_{1\leqslant i\leqslant \left( {{\begin{array}{*{20}c}
 r \hfill \\
 N \hfill \\
\end{array} }} \right)} {\left| {z_i } \right|^2} }\mbox{; } \\
 \\
 C_{ij} =\left. {\frac{\partial ^2\ln S}{\partial \bar {{z}'}_i \partial z_j
}} \right|_{{z}'=z} =\left. {\frac{1}{S}\left\{ {\frac{\partial
^2S}{\partial \bar {{z}'}_i \partial z_j }-\frac{1}{S}\frac{\partial
S}{\partial \bar {{z}'}_i }\frac{\partial S}{\partial z_j }} \right\}}
\right|_{{z}'=z} =\frac{1}{S}\left\{ {\delta _{ij} -\frac{1}{S}z_i \bar
{z}_j } \right\} \\
 \Rightarrow {\rm {\bf C}}=\frac{1}{S}\left\{ {{\rm {\bf I}}-{\rm {\bf
P}}_{\rm {\bf z}} } \right\}=\frac{1}{S}{\rm {\bf P}}_{\rm {\bf z}}^\bot
,\mbox{ } \\
 \mbox{where }{\rm {\bf P}}_{\rm {\bf z}} \mbox{ is the orthogonal projector
onto the vector }{\rm {\bf z}}\in C^{\left( {{\begin{array}{*{20}c}
 \mbox{r} \hfill \\
 \mbox{N} \hfill \\
\end{array} }} \right)} \\
 \end{array}
\]
\[
\begin{array}{c}
 \mbox{noting that }\left( {{\begin{array}{*{20}c}
 {\frac{\partial E}{\partial {\rm {\bf \bar {z}}}}} \hfill \\
 {\frac{\partial E}{\partial {\rm {\bf z}}}} \hfill \\
\end{array} }} \right)=\frac{1}{S}\left( {{\begin{array}{*{20}c}
 {\frac{\partial H}{\partial {\rm {\bf z}}}-E\frac{\partial S}{\partial {\rm
{\bf z}}}} \hfill \\
 {\frac{\partial H}{\partial {\rm {\bf \bar {z}}}}-E\frac{\partial
S}{\partial {\rm {\bf \bar {z}}}}} \hfill \\
\end{array} }} \right)=\frac{1}{S}\left( {{\begin{array}{*{20}c}
 {{\rm {\bf H}}\bullet {\rm {\bf z}}-E{\rm {\bf z}}} \hfill \\
 {{\rm {\bf \bar {H}}}\bullet {\rm {\bf \bar {z}}}-E{\rm {\bf \bar {z}}}}
\hfill \\
\end{array} }} \right)=\frac{1}{S}\left( {{\begin{array}{*{20}c}
 {\rm {\bf 0}} \hfill & {{\rm {\bf H}}-E{\rm {\bf I}}} \hfill \\
 {{\rm {\bf \bar {H}}}-E{\rm {\bf I}}} \hfill & {\rm {\bf 0}} \hfill \\
\end{array} }} \right)\left( {{\begin{array}{*{20}c}
 {{\rm {\bf \bar {z}}}} \hfill \\
 {\rm {\bf z}} \hfill \\
\end{array} }} \right) \\
 \mbox{the classical equations of motion }{\rm {\bf \dot {x}}}=\left\{ {{\rm
{\bf x}},E} \right\}\mbox{, }{\rm {\bf x}}={\rm {\bf z}},{\rm {\bf \bar
{z}}}\mbox{ }\mbox{give } \\
 \left( {{\begin{array}{*{20}c}
 {\rm {\bf 0}} \hfill & {\tfrac{i}{\hbar }{\rm {\bf C}}} \hfill \\
 {-\tfrac{i}{\hbar }{\rm {\bf \bar {C}}}} \hfill & {\rm {\bf 0}} \hfill \\
\end{array} }} \right)\left( {{\begin{array}{*{20}c}
 {{\rm {\bf \dot {\bar {z}}}}} \hfill \\
 {{\rm {\bf \dot {z}}}} \hfill \\
\end{array} }} \right)=\left( {{\begin{array}{*{20}c}
 {\frac{\partial E}{\partial {\rm {\bf \bar {z}}}}} \hfill \\
 {\frac{\partial E}{\partial {\rm {\bf z}}}} \hfill \\
\end{array} }} \right) \\
 \Rightarrow \left( {{\begin{array}{*{20}c}
 {\rm {\bf 0}} \hfill & {\tfrac{i}{\hbar }{\rm {\bf P}}_{\rm {\bf z}}^\bot }
\hfill \\
 {-\tfrac{i}{\hbar }{\rm {\bf P}}_{{\rm {\bf \bar {z}}}}^\bot } \hfill &
{\rm {\bf 0}} \hfill \\
\end{array} }} \right)\left( {{\begin{array}{*{20}c}
 {{\rm {\bf \dot {\bar {z}}}}} \hfill \\
 {{\rm {\bf \dot {z}}}} \hfill \\
\end{array} }} \right)=\left( {{\begin{array}{*{20}c}
 {\rm {\bf 0}} \hfill & {{\rm {\bf H}}-E{\rm {\bf I}}} \hfill \\
 {{\rm {\bf \bar {H}}}-E{\rm {\bf I}}} \hfill & {\rm {\bf 0}} \hfill \\
\end{array} }} \right)\left( {{\begin{array}{*{20}c}
 {{\rm {\bf \bar {z}}}} \hfill \\
 {\rm {\bf z}} \hfill \\
\end{array} }} \right) \\
 \mbox{as }\left( {{\rm {\bf H}}-E{\rm {\bf I}}} \right){\rm {\bf z}}\mbox{
is also }\bot \mbox{ }{\rm {\bf z}}\mbox{ this } \\
 \Rightarrow \left( {{\begin{array}{*{20}c}
 {-\tfrac{i}{\hbar }{\rm {\bf \dot {\bar {z}}}}} \hfill \\
 {\tfrac{i}{\hbar }{\rm {\bf \dot {z}}}} \hfill \\
\end{array} }} \right)=\left( {{\begin{array}{*{20}c}
 {\rm {\bf 0}} \hfill & {\rm {\bf I}} \hfill \\
 {-{\rm {\bf I}}} \hfill & {\rm {\bf 0}} \hfill \\
\end{array} }} \right)\left( {{\begin{array}{*{20}c}
 {\rm {\bf 0}} \hfill & {{\rm {\bf H}}-E{\rm {\bf I}}} \hfill \\
 {{\rm {\bf \bar {H}}}-E{\rm {\bf I}}} \hfill & {\rm {\bf 0}} \hfill \\
\end{array} }} \right)\left( {{\begin{array}{*{20}c}
 {{\rm {\bf \bar {z}}}} \hfill \\
 {\rm {\bf z}} \hfill \\
\end{array} }} \right)=\left( {{\begin{array}{*{20}c}
 {{\rm {\bf \bar {H}}}-E{\rm {\bf I}}} \hfill & {\rm {\bf 0}} \hfill \\
 {\rm {\bf 0}} \hfill & {{\rm {\bf H}}-E{\rm {\bf I}}} \hfill \\
\end{array} }} \right)\left( {{\begin{array}{*{20}c}
 {{\rm {\bf \bar {z}}}} \hfill \\
 {\rm {\bf z}} \hfill \\
\end{array} }} \right);\mbox{ } \\
 \mbox{The critical points of }E\left( {{\rm {\bf z}},{\rm {\bf \bar {z}}}}
\right)\mbox{ correspond to the stationary states of }H\mbox{ as} \\
 \left( {{\begin{array}{*{20}c}
 {\frac{\partial E}{\partial {\rm {\bf \bar {z}}}}} \hfill \\
 {\frac{\partial E}{\partial {\rm {\bf z}}}} \hfill \\
\end{array} }} \right)=\left( {{\begin{array}{*{20}c}
 {\rm {\bf 0}} \hfill \\
 {\rm {\bf 0}} \hfill \\
\end{array} }} \right)=\frac{1}{S}\left( {{\begin{array}{*{20}c}
 {\rm {\bf 0}} \hfill & {{\rm {\bf H}}-E{\rm {\bf I}}} \hfill \\
 {{\rm {\bf \bar {H}}}-E{\rm {\bf I}}} \hfill & {\rm {\bf 0}} \hfill \\
\end{array} }} \right)\left( {{\begin{array}{*{20}c}
 {{\rm {\bf \bar {z}}}} \hfill \\
 {\rm {\bf z}} \hfill \\
\end{array} }} \right)\Rightarrow \left| {\rm {\bf z}} \right\rangle \mbox{
are eigenstates of }H\mbox{ } \\
 \end{array}
\]
\[
\begin{array}{c}
 \mbox{noting that }\left( {{\begin{array}{*{20}c}
 {\frac{\partial E}{\partial {\rm {\bf \bar {z}}}}} \hfill \\
 {\frac{\partial E}{\partial {\rm {\bf z}}}} \hfill \\
\end{array} }} \right)=\frac{1}{S}\left( {{\begin{array}{*{20}c}
 {\frac{\partial H}{\partial {\rm {\bf z}}}-E\frac{\partial S}{\partial {\rm
{\bf z}}}} \hfill \\
 {\frac{\partial H}{\partial {\rm {\bf \bar {z}}}}-E\frac{\partial
S}{\partial {\rm {\bf \bar {z}}}}} \hfill \\
\end{array} }} \right)=\frac{1}{S}\left( {{\begin{array}{*{20}c}
 {{\rm {\bf H}}\bullet {\rm {\bf z}}-E{\rm {\bf z}}} \hfill \\
 {{\rm {\bf \bar {H}}}\bullet {\rm {\bf \bar {z}}}-E{\rm {\bf \bar {z}}}}
\hfill \\
\end{array} }} \right)=\frac{1}{S}\left( {{\begin{array}{*{20}c}
 {\rm {\bf 0}} \hfill & {{\rm {\bf H}}-E{\rm {\bf I}}} \hfill \\
 {{\rm {\bf \bar {H}}}-E{\rm {\bf I}}} \hfill & {\rm {\bf 0}} \hfill \\
\end{array} }} \right)\left( {{\begin{array}{*{20}c}
 {{\rm {\bf \bar {z}}}} \hfill \\
 {\rm {\bf z}} \hfill \\
\end{array} }} \right)\mbox{ } \\
 \mbox{the classical equations of motion }{\rm {\bf \dot {x}}}=\left\{ {{\rm
{\bf x}},E} \right\}\mbox{, }{\rm {\bf x}}={\rm {\bf z}},{\rm {\bf \bar
{z}}}\mbox{ }\mbox{give } \\
 \left( {{\begin{array}{*{20}c}
 {\rm {\bf 0}} \hfill & {\tfrac{i}{\hbar }{\rm {\bf C}}} \hfill \\
 {-\tfrac{i}{\hbar }{\rm {\bf \bar {C}}}} \hfill & {\rm {\bf 0}} \hfill \\
\end{array} }} \right)\left( {{\begin{array}{*{20}c}
 {{\rm {\bf \dot {\bar {z}}}}} \hfill \\
 {{\rm {\bf \dot {z}}}} \hfill \\
\end{array} }} \right)=\left( {{\begin{array}{*{20}c}
 {\frac{\partial E}{\partial {\rm {\bf \bar {z}}}}} \hfill \\
 {\frac{\partial E}{\partial {\rm {\bf z}}}} \hfill \\
\end{array} }} \right)\Rightarrow \left( {{\begin{array}{*{20}c}
 {\rm {\bf 0}} \hfill & {\tfrac{i}{\hbar }{\rm {\bf P}}_{\rm {\bf z}}^\bot }
\hfill \\
 {-\tfrac{i}{\hbar }{\rm {\bf P}}_{{\rm {\bf \bar {z}}}}^\bot } \hfill &
{\rm {\bf 0}} \hfill \\
\end{array} }} \right)\left( {{\begin{array}{*{20}c}
 {{\rm {\bf \dot {\bar {z}}}}} \hfill \\
 {{\rm {\bf \dot {z}}}} \hfill \\
\end{array} }} \right)=\left( {{\begin{array}{*{20}c}
 {\rm {\bf 0}} \hfill & {{\rm {\bf H}}-E{\rm {\bf I}}} \hfill \\
 {{\rm {\bf \bar {H}}}-E{\rm {\bf I}}} \hfill & {\rm {\bf 0}} \hfill \\
\end{array} }} \right)\left( {{\begin{array}{*{20}c}
 {{\rm {\bf \bar {z}}}} \hfill \\
 {\rm {\bf z}} \hfill \\
\end{array} }} \right) \\
 \mbox{as }\left( {{\rm {\bf H}}-E{\rm {\bf I}}} \right){\rm {\bf z}}\mbox{
is also }\bot \mbox{ }{\rm {\bf z}}\mbox{ this } \\
 \Rightarrow \left( {{\begin{array}{*{20}c}
 {-\tfrac{i}{\hbar }{\rm {\bf \dot {\bar {z}}}}} \hfill \\
 {\tfrac{i}{\hbar }{\rm {\bf \dot {z}}}} \hfill \\
\end{array} }} \right)=\left( {{\begin{array}{*{20}c}
 {\rm {\bf 0}} \hfill & {\rm {\bf I}} \hfill \\
 {-{\rm {\bf I}}} \hfill & {\rm {\bf 0}} \hfill \\
\end{array} }} \right)\left( {{\begin{array}{*{20}c}
 {\rm {\bf 0}} \hfill & {{\rm {\bf H}}-E{\rm {\bf I}}} \hfill \\
 {{\rm {\bf \bar {H}}}-E{\rm {\bf I}}} \hfill & {\rm {\bf 0}} \hfill \\
\end{array} }} \right)\left( {{\begin{array}{*{20}c}
 {{\rm {\bf \bar {z}}}} \hfill \\
 {\rm {\bf z}} \hfill \\
\end{array} }} \right)=\left( {{\begin{array}{*{20}c}
 {{\rm {\bf \bar {H}}}-E{\rm {\bf I}}} \hfill & {\rm {\bf 0}} \hfill \\
 {\rm {\bf 0}} \hfill & {{\rm {\bf H}}-E{\rm {\bf I}}} \hfill \\
\end{array} }} \right)\left( {{\begin{array}{*{20}c}
 {{\rm {\bf \bar {z}}}} \hfill \\
 {\rm {\bf z}} \hfill \\
\end{array} }} \right);\mbox{ } \\
 \mbox{The critical points of }E\left( {{\rm {\bf z}},{\rm {\bf \bar {z}}}}
\right)\mbox{ correspond to the stationary states of }H\mbox{ as } \\
 \left( {{\begin{array}{*{20}c}
 {\frac{\partial E}{\partial {\rm {\bf \bar {z}}}}} \hfill \\
 {\frac{\partial E}{\partial {\rm {\bf z}}}} \hfill \\
\end{array} }} \right)=\left( {{\begin{array}{*{20}c}
 {\rm {\bf 0}} \hfill \\
 {\rm {\bf 0}} \hfill \\
\end{array} }} \right)=\frac{1}{S}\left( {{\begin{array}{*{20}c}
 {\rm {\bf 0}} \hfill & {{\rm {\bf H}}-E{\rm {\bf I}}} \hfill \\
 {{\rm {\bf \bar {H}}}-E{\rm {\bf I}}} \hfill & {\rm {\bf 0}} \hfill \\
\end{array} }} \right)\left( {{\begin{array}{*{20}c}
 {{\rm {\bf \bar {z}}}} \hfill \\
 {\rm {\bf z}} \hfill \\
\end{array} }} \right)\mbox{ which}\Rightarrow \left| {\rm {\bf z}}
\right\rangle \mbox{ are eigenstates of }H\mbox{ } \\
 \mbox{ } \\
 \end{array}
\]
\[
\begin{array}{c}
 \mbox{We now consider the hamiltonian of the molecular system interacting
with the radiation } \\
 \mbox{field. If the point }\left( {{\rm {\bf z}}_0 ,{\rm {\bf \bar {z}}}_0
} \right)\mbox{ is a critical point of the molecular system, i.e }\left(
{{\rm {\bf H}}_{\rm {\bf 0}} -E{\rm {\bf I}}} \right)\bullet {\rm {\bf z}}_0
={\rm {\bf 0}}\mbox{ } \\
 \mbox{the Laplace transform of the linear response function becomes } \\
 R\left( z \right)=\left( {\left( {{\begin{array}{*{20}c}
 {\rm {\bf 0}} \hfill & {-z{\rm {\bf C}}} \hfill \\
 {z{\rm {\bf \bar {C}}}} \hfill & {\rm {\bf 0}} \hfill \\
\end{array} }} \right)-{\rm {\bf D}}^2E_0 \left( {{\rm {\bf z}}_0 ,{\rm {\bf
\bar {z}}}_0 } \right)} \right)^{-1}=\left( {{\begin{array}{*{20}c}
 {\rm {\bf 0}} \hfill & {-\left( {z{\rm {\bf I}}+\left( {{\rm {\bf \bar
{H}}}_{\rm {\bf 0}} -E{\rm {\bf I}}} \right)} \right)^{-1}} \hfill \\
 {z{\rm {\bf I}}-\left( {{\rm {\bf H}}_{\rm {\bf 0}} -E{\rm {\bf I}}}
\right)^{-1}} \hfill & {\rm {\bf 0}} \hfill \\
\end{array} }} \right) \\
 \mbox{where it }{\rm {\bf should}}\mbox{ }{\rm {\bf be}}\mbox{ }{\rm {\bf
observed}}\mbox{ that in this case the Hessian matrix is now block skew
diagonal } \\
 {\rm {\bf D}}^{\rm {\bf 2}}E\left( {{\rm {\bf z}}_0 ,{\rm {\bf \bar {z}}}_0
} \right)=\left( {{\begin{array}{*{20}c}
 {\frac{\partial ^2E_0 \left( {{\rm {\bf z}}_0 ,{\rm {\bf \bar {z}}}_0 }
\right)}{\partial {\rm {\bf \bar {z}}}\partial {\rm {\bf \bar {z}}}}} \hfill
& {\frac{\partial ^2E_0 \left( {{\rm {\bf z}}_0 ,{\rm {\bf \bar {z}}}_0 }
\right)}{\partial {\rm {\bf \bar {z}}}\partial {\rm {\bf z}}}} \hfill \\
 {\frac{\partial ^2E_0 \left( {{\rm {\bf z}}_0 ,{\rm {\bf \bar {z}}}_0 }
\right)}{\partial {\rm {\bf z}}\partial {\rm {\bf \bar {z}}}}} \hfill &
{\frac{\partial ^2E_0 \left( {{\rm {\bf z}}_0 ,{\rm {\bf \bar {z}}}_0 }
\right)}{\partial {\rm {\bf z}}\partial {\rm {\bf z}}}} \hfill \\
\end{array} }} \right)=\left( {{\begin{array}{*{20}c}
 {\rm {\bf 0}} \hfill & {\frac{\partial ^2E_0 \left( {{\rm {\bf z}}_0 ,{\rm
{\bf \bar {z}}}_0 } \right)}{\partial {\rm {\bf \bar {z}}}\partial {\rm {\bf
z}}}} \hfill \\
 {\frac{\partial ^2E_0 \left( {{\rm {\bf z}}_0 ,{\rm {\bf \bar {z}}}_0 }
\right)}{\partial {\rm {\bf z}}\partial {\rm {\bf \bar {z}}}}} \hfill & {\rm
{\bf 0}} \hfill \\
\end{array} }} \right). \\
 \mbox{The linear response of the 1-particle reduced density matrix in this
case is thus } \\
 \delta ^1\rho ^1\left( {{\rm {\bf z}}_0 ,{\rm {\bf \bar {z}}}_0 }
\right)_{ij} \left( z \right)=\left( {{\begin{array}{*{20}c}
 {\frac{\partial \rho ^1\left( {{\rm {\bf z}}_0 ,{\rm {\bf \bar {z}}}_0 }
\right)_{ij} }{\partial {\rm {\bf \bar {z}}}}} \hfill & {\frac{\partial \rho
^1\left( {{\rm {\bf z}}_0 ,{\rm {\bf \bar {z}}}_0 } \right)_{ij} }{\partial
{\rm {\bf z}}}} \hfill \\
\end{array} }} \right)\left( {{\begin{array}{*{20}c}
 {\rm {\bf 0}} \hfill & {-\left( {z{\rm {\bf I}}+\left( {{\rm {\bf \bar
{H}}}_{\rm {\bf 0}} -E{\rm {\bf I}}} \right)} \right)^{-1}} \hfill \\
 {z{\rm {\bf I}}-\left( {{\rm {\bf H}}_{\rm {\bf 0}} -E{\rm {\bf I}}}
\right)^{-1}} \hfill & {\rm {\bf 0}} \hfill \\
\end{array} }} \right){\rm {\bf D}}^{\rm {\bf 1}}F\left( {{\rm {\bf z}}_0
,{\rm {\bf \bar {z}}}_0 ,z} \right) \\
 \mbox{where }{\rm {\bf D}}^{\rm {\bf 1}}F\left( {{\rm {\bf z}}_0 ,{\rm {\bf
\bar {z}}}_0 ,z} \right)=\left( {{\begin{array}{*{20}c}
 {{\rm {\bf F}}\left( z \right){\rm {\bf z}}_0 } \hfill \\
 {{\rm {\bf \bar {F}}}\left( z \right){\rm {\bf \bar {z}}}_0 } \hfill \\
\end{array} }} \right)=\left( {{\begin{array}{*{20}c}
 {{\rm {\bf F}}\left( z \right)} \hfill & {\rm {\bf 0}} \hfill \\
 {\rm {\bf 0}} \hfill & {{\rm {\bf \bar {F}}}\left( z \right)} \hfill \\
\end{array} }} \right)\left( {{\begin{array}{*{20}c}
 {{\rm {\bf z}}_0 } \hfill \\
 {{\rm {\bf \bar {z}}}_0 } \hfill \\
\end{array} }} \right)\mbox{; } \\
 \end{array}
\]
\[
\begin{array}{c}
 \mbox{The hamiltonian of the molecular system plus perturbation is }H=H_0
+F\left( t \right)\mbox{ } \\
 \mbox{has the matrix representation wrt to the complete orthormal basis
}\left\{ {\Phi _\mbox{i} } \right\}\equiv \underline{\Phi }\mbox{ } \\
 {\rm {\bf H}}=\left\langle {\underline{\Phi }} \mathrel{\left| {\vphantom
{{\underline{\Phi }} {H\underline{\Phi }}}} \right.
\kern-\nulldelimiterspace} {H\underline{\Phi }} \right\rangle ={\rm {\bf
H}}_{\rm {\bf 0}} +{\rm {\bf F}}\left( t \right).\mbox{ The derrivatives of
the energy can be expressed in terms } \\
 \mbox{of these matrices as } \\
 \frac{\partial E\left( {{\rm {\bf z}}_0 ,{\rm {\bf \bar {z}}}_0 }
\right)}{\partial {\rm {\bf z}}}=\frac{1}{S}\left\langle {{\rm {\bf z}}_0 }
\mathrel{\left| {\vphantom {{{\rm {\bf z}}_0 } {\left( {H_0 -EI+F\left( t
\right)} \right)\underline{\Phi }}}} \right. \kern-\nulldelimiterspace}
{\left( {H_0 -EI+F\left( t \right)} \right)\underline{\Phi }} \right\rangle
=\frac{1}{S}\left\{ {\left\langle {{\rm {\bf z}}_0 } \mathrel{\left|
{\vphantom {{{\rm {\bf z}}_0 } {\left( {H_0 -EI} \right)\underline{\Phi }}}}
\right. \kern-\nulldelimiterspace} {\left( {H_0 -EI} \right)\underline{\Phi
}} \right\rangle +\sum\limits_{ij} {f_{kl} \left( t \right)\mbox{ }}
\left\langle {{\rm {\bf z}}_0 } \mathrel{\left| {\vphantom {{{\rm {\bf z}}_0
} {\mbox{a}_k^\dag \left( t \right)\mbox{a}_l \left( t
\right)\underline{\Phi }}}} \right. \kern-\nulldelimiterspace}
{\mbox{a}_k^\dag \left( t \right)\mbox{a}_l \left( t \right)\underline{\Phi
}} \right\rangle } \right\} \\
 =\frac{1}{S}\left( {{\rm {\bf \bar {H}}}_{\rm {\bf 0}} -E{\rm {\bf I}}+{\rm
{\bf \bar {F}}}\left( t \right)} \right)\bullet {\rm {\bf \bar {z}}}_0 \\
 \frac{\partial ^2E\left( {{\rm {\bf z}}_0 ,{\rm {\bf \bar {z}}}_0 }
\right)}{\partial {\rm {\bf \bar {z}}}\partial {\rm {\bf
z}}}=\frac{1}{S}\left\langle {\underline{\Phi }} \mathrel{\left| {\vphantom
{{\underline{\Phi }} {\left( {H_0 -EI+F\left( t \right)}
\right)\underline{\Phi }}}} \right. \kern-\nulldelimiterspace} {\left( {H_0
-EI+F\left( t \right)} \right)\underline{\Phi }} \right\rangle
=\frac{1}{S}\left\{ {\left\langle {\underline{\Phi }} \mathrel{\left|
{\vphantom {{\underline{\Phi }} {\left( {H_0 -EI} \right)\underline{\Phi
}}}} \right. \kern-\nulldelimiterspace} {\left( {H_0 -EI}
\right)\underline{\Phi }} \right\rangle +\sum\limits_{ij} {f_{kl} \left( t
\right)\mbox{ }} \left\langle {\underline{\Phi }} \mathrel{\left| {\vphantom
{{\underline{\Phi }} {\mbox{a}_k^\dag \left( t \right)\mbox{a}_l \left( t
\right)\underline{\Phi }}}} \right. \kern-\nulldelimiterspace}
{\mbox{a}_k^\dag \left( t \right)\mbox{a}_l \left( t \right)\underline{\Phi
}} \right\rangle } \right\} \\
 =\frac{1}{S}\left( {{\rm {\bf \bar {H}}}_{\rm {\bf 0}} -E{\rm {\bf I}}+{\rm
{\bf \bar {F}}}\left( t \right)} \right) \\
 \mbox{Where }E=E\left( {{\rm {\bf z}}_0 ,{\rm {\bf \bar {z}}}_0 }
\right)\mbox{. Expanding the energy about the point }\left( {{\rm {\bf z}}_0
,{\rm {\bf \bar {z}}}_0 } \right)\mbox{ gives the } \\
 exact\mbox{ expression } \\
 E\left( {{\rm {\bf z}}_0 +\delta {\rm {\bf z}},{\rm {\bf \bar {z}}}_0
+\delta {\rm {\bf \bar {z}}}} \right)=E\left( {{\rm {\bf z}}_0 ,{\rm {\bf
\bar {z}}}_0 } \right)+\delta {\rm {\bf z}}^\dag \bullet \frac{\partial
E\left( {{\rm {\bf z}}_0 ,{\rm {\bf \bar {z}}}_0 } \right)}{\partial {\rm
{\bf \bar {z}}}}+\frac{\partial E\left( {{\rm {\bf z}}_0 ,{\rm {\bf \bar
{z}}}_0 } \right)}{\partial {\rm {\bf z}}}\bullet \delta {\rm {\bf z}} \\
 +\frac{1}{2}\left\{ {\delta {\rm {\bf z}}^\dag \bullet \frac{\partial
^2E\left( {{\rm {\bf z}}_0 ,{\rm {\bf \bar {z}}}_0 } \right)}{\partial {\rm
{\bf \bar {z}}}\partial {\rm {\bf z}}}\bullet \delta {\rm {\bf z}}+\delta
{\rm {\bf \bar {z}}}^\dag \bullet \frac{\partial ^2E\left( {{\rm {\bf z}}_0
,{\rm {\bf \bar {z}}}_0 } \right)}{\partial {\rm {\bf z}}\partial {\rm {\bf
\bar {z}}}}\bullet \delta {\rm {\bf \bar {z}}}} \right\} \\
 \mbox{ Linear response to the external perturbation is obtain by
approximating the } \\
 \mbox{energy function around the point }\left( {{\rm {\bf z}}_0 ,{\rm {\bf
\bar {z}}}_0 } \right)\mbox{ by only retaining terms up to first order } \\
 \mbox{in the perturbation } \\
 E\left( {{\rm {\bf z}}_0 +\delta {\rm {\bf z}},{\rm {\bf \bar {z}}}_0
+\delta {\rm {\bf \bar {z}}}} \right)\approx \tilde {E}\left( {{\rm {\bf
z}}_0 +\delta {\rm {\bf z}},{\rm {\bf \bar {z}}}_0 +\delta {\rm {\bf \bar
{z}}}} \right)=E\left( {{\rm {\bf z}}_0 ,{\rm {\bf \bar {z}}}_0 }
\right)+\delta {\rm {\bf z}}^\dag \bullet \left\langle {\underline{\Phi }}
\mathrel{\left| {\vphantom {{\underline{\Phi }} {\left( {H_0 +F\left( t
\right)} \right){\rm {\bf z}}_0 }}} \right. \kern-\nulldelimiterspace}
{\left( {H_0 +F\left( t \right)} \right){\rm {\bf z}}_0 } \right\rangle \\
 +\left\langle {{\rm {\bf z}}_0 } \mathrel{\left| {\vphantom {{{\rm {\bf
z}}_0 } {\left( {H_0 +F\left( t \right)} \right)\underline{\Phi }}}} \right.
\kern-\nulldelimiterspace} {\left( {H_0 +F\left( t \right)}
\right)\underline{\Phi }} \right\rangle \bullet \delta {\rm {\bf z}}+\delta
{\rm {\bf z}}^\dag \bullet \left\langle {\underline{\Phi }} \mathrel{\left|
{\vphantom {{\underline{\Phi }} {H_0 \underline{\Phi }}}} \right.
\kern-\nulldelimiterspace} {H_0 \underline{\Phi }} \right\rangle \bullet
\delta {\rm {\bf z}} \\
 \mbox{to give the local linear response energy gradient } \\
 \left( {{\begin{array}{*{20}c}
 {\frac{\partial \tilde {E}\left( {{\rm {\bf z}}_0 ,{\rm {\bf \bar {z}}}_0 }
\right)}{\partial \delta {\rm {\bf \bar {z}}}}} \hfill \\
 {\frac{\partial \tilde {E}\left( {{\rm {\bf z}}_0 ,{\rm {\bf \bar {z}}}_0 }
\right)}{\partial \delta {\rm {\bf z}}}} \hfill \\
\end{array} }} \right)=\frac{1}{S}\left( {{\begin{array}{*{20}c}
 {{\rm {\bf H}}_{\rm {\bf 0}} -E{\rm {\bf I}}+{\rm {\bf F}}\left( t \right)}
\hfill & {\rm {\bf 0}} \hfill \\
 {\rm {\bf 0}} \hfill & {{\rm {\bf \bar {H}}}_{\rm {\bf 0}} -E{\rm {\bf
I}}+{\rm {\bf \bar {F}}}\left( t \right)} \hfill \\
\end{array} }} \right)\left( {{\begin{array}{*{20}c}
 {{\rm {\bf z}}_0 } \hfill \\
 {{\rm {\bf \bar {z}}}_0 } \hfill \\
\end{array} }} \right) \\
 +\frac{1}{S}\left( {{\begin{array}{*{20}c}
 {{\rm {\bf H}}_{\rm {\bf 0}} -E{\rm {\bf I}}} \hfill & {\rm {\bf 0}} \hfill
\\
 {\rm {\bf 0}} \hfill & {{\rm {\bf \bar {H}}}_{\rm {\bf 0}} -E{\rm {\bf I}}}
\hfill \\
\end{array} }} \right)\left( {{\begin{array}{*{20}c}
 {\delta {\rm {\bf z}}} \hfill \\
 {\delta {\rm {\bf \bar {z}}}} \hfill \\
\end{array} }} \right) \\
 \mbox{ } \\
 \end{array}
\begin{array}{c}
 \Rightarrow \delta ^1\rho ^1\left( {{\rm {\bf z}}_0 ,{\rm {\bf \bar {z}}}_0
} \right)_{ij} \left( z \right)=\left( {{\begin{array}{*{20}c}
 {{\rm {\bf \bar {z}}}_0^\dag \bullet \overline {\left\langle
{\underline{\Phi }} \mathrel{\left| {\vphantom {{\underline{\Phi }}
{\mbox{a}_j^\dag \mbox{a}_i \underline{\Phi }}}} \right.
\kern-\nulldelimiterspace} {\mbox{a}_j^\dag \mbox{a}_i \underline{\Phi }}
\right\rangle } } \hfill & {{\rm {\bf z}}_0^\dag \bullet \left\langle
{\underline{\Phi }} \mathrel{\left| {\vphantom {{\underline{\Phi }}
{\mbox{a}_j^\dag \mbox{a}_i \underline{\Phi }}}} \right.
\kern-\nulldelimiterspace} {\mbox{a}_j^\dag \mbox{a}_i \underline{\Phi }}
\right\rangle } \hfill \\
\end{array} }} \right)\left( {{\begin{array}{*{20}c}
 {\rm {\bf 0}} \hfill & {\rm {\bf 1}} \hfill \\
 {\rm {\bf 1}} \hfill & {\rm {\bf 0}} \hfill \\
\end{array} }} \right)\left( {{\begin{array}{*{20}c}
 {\rm {\bf 0}} \hfill & {\rm {\bf 1}} \hfill \\
 {\rm {\bf 1}} \hfill & {\rm {\bf 0}} \hfill \\
\end{array} }} \right)\times \mbox{ } \\
 \mbox{ }\times \left( {{\begin{array}{*{20}c}
 {\rm {\bf 0}} \hfill & {-\left( {z{\rm {\bf I}}+\left( {{\rm {\bf \bar
{H}}}_{\rm {\bf 0}} -E{\rm {\bf I}}} \right)} \right)^{-1}} \hfill \\
 {z{\rm {\bf I}}-\left( {{\rm {\bf H}}_{\rm {\bf 0}} -E{\rm {\bf I}}}
\right)^{-1}} \hfill & {\rm {\bf 0}} \hfill \\
\end{array} }} \right)\left( {{\begin{array}{*{20}c}
 {{\rm {\bf F}}\left( z \right)} \hfill & {\rm {\bf 0}} \hfill \\
 {\rm {\bf 0}} \hfill & {{\rm {\bf \bar {F}}}\left( z \right)} \hfill \\
\end{array} }} \right)\left( {{\begin{array}{*{20}c}
 {{\rm {\bf z}}_0 } \hfill \\
 {{\rm {\bf \bar {z}}}_0 } \hfill \\
\end{array} }} \right) \\
 \\
 \end{array}
\]
\[
\begin{array}{c}
 =\left( {{\begin{array}{*{20}c}
 {\left\langle {\underline{{\rm {\bf z}}_0 }} \mathrel{\left| {\vphantom
{{\underline{{\rm {\bf z}}_0 }} {\mbox{a}_j^\dag \mbox{a}_i \underline{\Phi
}}}} \right. \kern-\nulldelimiterspace} {\mbox{a}_j^\dag \mbox{a}_i
\underline{\Phi }} \right\rangle } \hfill & {\overline {\left\langle
{\underline{{\rm {\bf z}}_0 }} \mathrel{\left| {\vphantom {{\underline{{\rm
{\bf z}}_0 }} {\mbox{a}_j^\dag \mbox{a}_i \underline{\Phi }}}} \right.
\kern-\nulldelimiterspace} {\mbox{a}_j^\dag \mbox{a}_i \underline{\Phi }}
\right\rangle } } \hfill \\
\end{array} }} \right)\times \mbox{ } \\
 \times \left( {{\begin{array}{*{20}c}
 {\left\langle {\underline{\Phi }} \mathrel{\left| {\vphantom
{{\underline{\Phi }} {\left\{ {zI-\left( {H_0 -EI} \right)^{-1}}
\right\}\underline{\Phi }}}} \right. \kern-\nulldelimiterspace} {\left\{
{zI-\left( {H_0 -EI} \right)^{-1}} \right\}\underline{\Phi }} \right\rangle
} \hfill & {\rm {\bf 0}} \hfill \\
 {\rm {\bf 0}} \hfill & {-\overline {\left\langle {\underline{\Phi }}
\mathrel{\left| {\vphantom {{\underline{\Phi }} {\left\{ {\left( {\bar
{z}I+\left( {H_0 -EI} \right)} \right)^{-1}} \right\}\underline{\Phi }}}}
\right. \kern-\nulldelimiterspace} {\left\{ {\left( {\bar {z}I+\left( {H_0
-EI} \right)} \right)^{-1}} \right\}\underline{\Phi }} \right\rangle } }
\hfill \\
\end{array} }} \right)\left( {{\begin{array}{*{20}c}
 {\left\langle {\underline{\Phi }} \mathrel{\left| {\vphantom
{{\underline{\Phi }} {\left( {{\rm {\bf a}}^{\rm {\bf \dag }}{\rm {\bf a}}}
\right)\underline{{\rm {\bf z}}_0 }}}} \right. \kern-\nulldelimiterspace}
{\left( {{\rm {\bf a}}^{\rm {\bf \dag }}{\rm {\bf a}}}
\right)\underline{{\rm {\bf z}}_0 }} \right\rangle } \hfill \\
 {\overline {\left\langle {\underline{\Phi }} \mathrel{\left| {\vphantom
{{\underline{\Phi }} {\left( {{\rm {\bf a}}^{\rm {\bf \dag }}{\rm {\bf a}}}
\right)\underline{{\rm {\bf z}}_0 }}}} \right. \kern-\nulldelimiterspace}
{\left( {{\rm {\bf a}}^{\rm {\bf \dag }}{\rm {\bf a}}}
\right)\underline{{\rm {\bf z}}_0 }} \right\rangle } } \hfill \\
\end{array} }} \right)\left( {{\begin{array}{*{20}c}
 {\begin{array}{l}
 {\rm {\bf f}}\left( z \right) \\
 \\
 \end{array}} \hfill \\
 {\overline {{\rm {\bf f}}\left( z \right)} } \hfill \\
\end{array} }} \right)= \\
 \mbox{ } \\
 \sum\limits_{1\leqslant k,l\leqslant r} {\left\langle {{\rm {\bf z}}_0 }
\mathrel{\left| {\vphantom {{{\rm {\bf z}}_0 } {\left( {\mbox{a}_j^\dag
\mbox{a}_i \left\{ {zI-\left( {H_0 -EI} \right)^{-1}}
\right\}\mbox{a}_k^\dag \mbox{a}_l } \right){\rm {\bf z}}_0 }}} \right.
\kern-\nulldelimiterspace} {\left( {\mbox{a}_j^\dag \mbox{a}_i \left\{
{zI-\left( {H_0 -EI} \right)^{-1}} \right\}\mbox{a}_k^\dag \mbox{a}_l }
\right){\rm {\bf z}}_0 } \right\rangle f_{kl} \left( z \right)-\left\langle
{{\rm {\bf z}}_0 } \mathrel{\left| {\vphantom {{{\rm {\bf z}}_0 } {\left(
{\mbox{a}_k^\dag \mbox{a}_l \left\{ {zI+\left( {H_0 -EI} \right)^{-1}}
\right\}\mbox{a}_j^\dag \mbox{a}_i } \right){\rm {\bf z}}_0 }}} \right.
\kern-\nulldelimiterspace} {\left( {\mbox{a}_k^\dag \mbox{a}_l \left\{
{zI+\left( {H_0 -EI} \right)^{-1}} \right\}\mbox{a}_j^\dag \mbox{a}_i }
\right){\rm {\bf z}}_0 } \right\rangle \overline {f_{lk} \left( z \right)} }
\\
 =\sum\limits_{kl} {R_{ijkl} \left( z \right)f_{kl} \left( z \right)} \\
 \mbox{ } \\
 \end{array}
\]
Where $R_{ijkl} \left( z \right)$ are precisely the exact linear response
functions obtained previously.

\underline {\textbf{4.4 Linear Response Described by
}}$\raise0.7ex\hbox{${\mbox{SU}\left( r \right)}$} \!\mathord{\left/
{\vphantom {{\mbox{SU}\left( r \right)} {\mbox{USp}\left( r
\right)}}}\right.\kern-\nulldelimiterspace}\!\lower0.7ex\hbox{${\mbox{USp}\left(
r \right)}$}$\underline {\textbf{-Coherent States .}}

We now use the manifold of AGP coherent states in which to perform the
stationary phase approximation and to consider approximate linear response
functions. Within the context of the classical equations of motion in the
associated phase the linear response of the first order reduced density
matrix produced by the radiation field is then given by
\[
\begin{array}{l}
 \delta ^1\rho ^1\left( {{\rm {\bf z}}_0 ,{\rm {\bf \bar {z}}}_0 }
\right)_{ij} \left( z \right)= \\
 \left( {{\begin{array}{*{20}c}
 {\frac{\partial \rho ^1\left( {{\rm {\bf z}}_0 ,{\rm {\bf \bar {z}}}_0 }
\right)_{ij} }{\partial {\rm {\bf \bar {z}}}}} \hfill & {\frac{\partial \rho
^1\left( {{\rm {\bf z}}_0 ,{\rm {\bf \bar {z}}}_0 } \right)_{ij} }{\partial
{\rm {\bf z}}}} \hfill \\
\end{array} }} \right)\left( {{\begin{array}{*{20}c}
 {\rm {\bf 0}} \hfill & {\rm {\bf 1}} \hfill \\
 {\rm {\bf 1}} \hfill & {\rm {\bf 0}} \hfill \\
\end{array} }} \right)\left( {{\begin{array}{*{20}c}
 {\rm {\bf 0}} \hfill & {\rm {\bf 1}} \hfill \\
 {\rm {\bf 1}} \hfill & {\rm {\bf 0}} \hfill \\
\end{array} }} \right)\left( {\left( {{\begin{array}{*{20}c}
 {\rm {\bf 0}} \hfill & {-z{\rm {\bf C}}} \hfill \\
 {z{\rm {\bf \bar {C}}}} \hfill & {\rm {\bf 0}} \hfill \\
\end{array} }} \right)-{\rm {\bf D}}^2E_0 \left( {{\rm {\bf z}}_0 ,{\rm {\bf
\bar {z}}}_0 } \right)} \right)^{-1}{\rm {\bf D}}^{\rm {\bf 1}}F\left( {{\rm
{\bf z}}_0 ,{\rm {\bf \bar {z}}}_0 ,z} \right) \\
 =\left( {{\begin{array}{*{20}c}
 {\frac{\partial \rho ^1\left( {{\rm {\bf z}}_0 ,{\rm {\bf \bar {z}}}_0 }
\right)_{ij} }{\partial {\rm {\bf z}}}} \hfill & {\frac{\partial \rho
^1\left( {{\rm {\bf z}}_0 ,{\rm {\bf \bar {z}}}_0 } \right)_{ij} }{\partial
{\rm {\bf \bar {z}}}}} \hfill \\
\end{array} }} \right)\left( {\left( {{\begin{array}{*{20}c}
 {z{\rm {\bf \bar {C}}}} \hfill & {\rm {\bf 0}} \hfill \\
 {\rm {\bf 0}} \hfill & {-z{\rm {\bf C}}} \hfill \\
\end{array} }} \right)-\left( {{\begin{array}{*{20}c}
 {\frac{\partial ^2E_0 \left( {{\rm {\bf z}}_0 ,{\rm {\bf \bar {z}}}_0 }
\right)}{\partial {\rm {\bf z}}\partial {\rm {\bf \bar {z}}}}} \hfill &
{\frac{\partial ^2E_0 \left( {{\rm {\bf z}}_0 ,{\rm {\bf \bar {z}}}_0 }
\right)}{\partial {\rm {\bf z}}\partial {\rm {\bf z}}}} \hfill \\
 {\frac{\partial ^2E_0 \left( {{\rm {\bf z}}_0 ,{\rm {\bf \bar {z}}}_0 }
\right)}{\partial {\rm {\bf \bar {z}}}\partial {\rm {\bf \bar {z}}}}} \hfill
& {\frac{\partial ^2E_0 \left( {{\rm {\bf z}}_0 ,{\rm {\bf \bar {z}}}_0 }
\right)}{\partial {\rm {\bf z}}\partial {\rm {\bf \bar {z}}}}} \hfill \\
\end{array} }} \right)} \right)^{-1}\left( {{\begin{array}{*{20}c}
 {\begin{array}{l}
 \frac{\partial F\left( {{\rm {\bf z}}_0 ,{\rm {\bf \bar {z}}}_0 ,z}
\right)}{\partial {\rm {\bf z}}} \\
 \\
 \end{array}} \hfill \\
 {\overline {\frac{\partial F\left( {{\rm {\bf z}}_0 ,{\rm {\bf \bar {z}}}_0
,z} \right)}{\partial {\rm {\bf z}}}} } \hfill \\
\end{array} }} \right) \\
 \frac{\partial F\left( {{\rm {\bf z}}_0 ,{\rm {\bf \bar {z}}}_0 ,z}
\right)}{\partial {\rm {\bf z}}}=\sum\limits_{kl} {\left. {\frac{\partial
\left\langle {{\rm {\bf z}}} \mathrel{\left| {\vphantom {{\rm {\bf z}}
{\mbox{a}_k^\dag \mbox{a}_l {\rm {\bf z}}}}} \right.
\kern-\nulldelimiterspace} {\mbox{a}_k^\dag \mbox{a}_l {\rm {\bf z}}}
\right\rangle }{\partial {\rm {\bf z}}}} \right|} _{{\rm {\bf z}}={\rm {\bf
z}}_{\rm {\bf 0}} } f_{kl} \left( z \right)\mbox{; and }\frac{\partial \rho
^1\left( {{\rm {\bf z}}_0 ,{\rm {\bf \bar {z}}}_0 } \right)_{ij} }{\partial
{\rm {\bf z}}}=\left. {\frac{\partial \left\langle {{\rm {\bf z}}}
\mathrel{\left| {\vphantom {{\rm {\bf z}} {\mbox{a}_j^\dag \mbox{a}_i {\rm
{\bf z}}}}} \right. \kern-\nulldelimiterspace} {\mbox{a}_j^\dag \mbox{a}_i
{\rm {\bf z}}} \right\rangle }{\partial {\rm {\bf z}}}} \right|_{{\rm {\bf
z}}={\rm {\bf z}}_{\rm {\bf 0}} } \\
 \Rightarrow \delta ^1\rho ^1\left( {{\rm {\bf z}}_0 ,{\rm {\bf \bar {z}}}_0
} \right)_{ij} \left( z \right)=\left( {{\begin{array}{*{20}c}
 {\left. {\frac{\partial \left\langle {{\rm {\bf z}}} \mathrel{\left|
{\vphantom {{\rm {\bf z}} {\mbox{a}_j^\dag \mbox{a}_i {\rm {\bf z}}}}}
\right. \kern-\nulldelimiterspace} {\mbox{a}_j^\dag \mbox{a}_i {\rm {\bf
z}}} \right\rangle }{\partial {\rm {\bf z}}}} \right|_{{\rm {\bf z}}={\rm
{\bf z}}_{\rm {\bf 0}} } } \hfill & {\left. {\frac{\partial \left\langle
{{\rm {\bf z}}} \mathrel{\left| {\vphantom {{\rm {\bf z}} {\mbox{a}_j^\dag
\mbox{a}_i {\rm {\bf z}}}}} \right. \kern-\nulldelimiterspace}
{\mbox{a}_j^\dag \mbox{a}_i {\rm {\bf z}}} \right\rangle }{\partial {\rm
{\bf \bar {z}}}}} \right|_{{\rm {\bf z}}={\rm {\bf z}}_{\rm {\bf 0}} } }
\hfill \\
\end{array} }} \right)\times \\
 \times \left( {\left( {{\begin{array}{*{20}c}
 {z{\rm {\bf \bar {C}}}} \hfill & {\rm {\bf 0}} \hfill \\
 {\rm {\bf 0}} \hfill & {-z{\rm {\bf C}}} \hfill \\
\end{array} }} \right)-\left( {{\begin{array}{*{20}c}
 {\frac{\partial ^2E_0 \left( {{\rm {\bf z}}_0 ,{\rm {\bf \bar {z}}}_0 }
\right)}{\partial {\rm {\bf z}}\partial {\rm {\bf \bar {z}}}}} \hfill &
{\frac{\partial ^2E_0 \left( {{\rm {\bf z}}_0 ,{\rm {\bf \bar {z}}}_0 }
\right)}{\partial {\rm {\bf z}}\partial {\rm {\bf z}}}} \hfill \\
 {\frac{\partial ^2E_0 \left( {{\rm {\bf z}}_0 ,{\rm {\bf \bar {z}}}_0 }
\right)}{\partial {\rm {\bf \bar {z}}}\partial {\rm {\bf \bar {z}}}}} \hfill
& {\frac{\partial ^2E_0 \left( {{\rm {\bf z}}_0 ,{\rm {\bf \bar {z}}}_0 }
\right)}{\partial {\rm {\bf z}}\partial {\rm {\bf \bar {z}}}}} \hfill \\
\end{array} }} \right)} \right)^{-1}\left( {{\begin{array}{*{20}c}
 {\left. {\left. {\frac{\partial \left\langle {{\rm {\bf z}}}
\mathrel{\left| {\vphantom {{\rm {\bf z}} {\left( {{\rm {\bf a}}^{\rm {\bf
\dag }}{\rm {\bf a}}} \right){\rm {\bf z}}}}} \right.
\kern-\nulldelimiterspace} {\left( {{\rm {\bf a}}^{\rm {\bf \dag }}{\rm {\bf
a}}} \right){\rm {\bf z}}} \right\rangle }{\partial {\rm {\bf z}}}} \right|}
\right|_{{\rm {\bf z}}={\rm {\bf z}}_{\rm {\bf 0}} } } \hfill & {\rm {\bf
0}} \hfill \\
 {\rm {\bf 0}} \hfill & {\overline {\left. {\left. {\frac{\partial
\left\langle {{\rm {\bf z}}} \mathrel{\left| {\vphantom {{\rm {\bf z}}
{\left( {{\rm {\bf a}}^{\rm {\bf \dag }}{\rm {\bf a}}} \right){\rm {\bf
z}}}}} \right. \kern-\nulldelimiterspace} {\left( {{\rm {\bf a}}^{\rm {\bf
\dag }}{\rm {\bf a}}} \right){\rm {\bf z}}} \right\rangle }{\partial {\rm
{\bf z}}}} \right|} \right|_{{\rm {\bf z}}={\rm {\bf z}}_{\rm {\bf 0}} } } }
\hfill \\
\end{array} }} \right)\left( {{\begin{array}{*{20}c}
 {\begin{array}{l}
 {\rm {\bf f}}\left( z \right) \\
 \\
 \end{array}} \hfill \\
 {\overline {{\rm {\bf f}}\left( z \right)} } \hfill \\
\end{array} }} \right) \\
 \end{array}
\]
In general the off diagonal blocks $\left\{ {\frac{\partial ^2E_0 \left(
{{\rm {\bf z}}_0 ,{\rm {\bf \bar {z}}}_0 } \right)}{\partial {\rm {\bf
z}}\partial {\rm {\bf z}}},\frac{\partial ^2E_0 \left( {{\rm {\bf z}}_0
,{\rm {\bf \bar {z}}}_0 } \right)}{\partial {\rm {\bf \bar {z}}}\partial
{\rm {\bf \bar {z}}}}} \right\}$\underline { } are not zero
\[
\begin{array}{c}
 H_0 \left( {{\rm {\bf z}},{\rm {\bf \bar {z}}}} \right)=\left\langle
{\mbox{exp}\sum\limits_{\begin{array}{l}
 -s\leqslant i<j\leqslant s \\
 \left| i \right|\ne \left| j \right| \\
 \end{array}} {\mbox{z}_\mbox{+} _{ij} q_{ij}^\dag } \mbox{
exp}\sum\limits_{\begin{array}{l}
 \mbox{-}s\leqslant i\leqslant s \\
 \left| i \right|=\left| j \right| \\
 \end{array}} {\left\{ {\lambda _{ij} u_{ij}^ } \right\}}
g_\mbox{0}^{\tfrac{\mbox{N}}{\mbox{2}}} } \mathrel{\left| {\vphantom
{{\mbox{exp}\sum\limits_{\begin{array}{l}
 -s\leqslant i<j\leqslant s \\
 \left| i \right|\ne \left| j \right| \\
 \end{array}} {\mbox{z}_\mbox{+} _{ij} q_{ij}^\dag } \mbox{
exp}\sum\limits_{\begin{array}{l}
 \mbox{-}s\leqslant i\leqslant s \\
 \left| i \right|=\left| j \right| \\
 \end{array}} {\left\{ {\lambda _{ij} u_{ij}^ } \right\}}
g_\mbox{0}^{\tfrac{\mbox{N}}{\mbox{2}}} } {H_0
\mbox{exp}\sum\limits_{\begin{array}{l}
 -s\leqslant i<j\leqslant s \\
 \left| i \right|\ne \left| j \right| \\
 \end{array}} {\mbox{z}_\mbox{+} _{ij} q_{ij}^\dag } \mbox{
exp}\sum\limits_{\begin{array}{l}
 \mbox{-}s\leqslant i\leqslant s \\
 \left| i \right|=\left| j \right| \\
 \end{array}} {\left\{ {\lambda _{ij} u_{ij}^ } \right\}}
g_\mbox{0}^{\tfrac{\mbox{N}}{\mbox{2}}} }}} \right.
\kern-\nulldelimiterspace} {H_0 \mbox{exp}\sum\limits_{\begin{array}{l}
 -s\leqslant i<j\leqslant s \\
 \left| i \right|\ne \left| j \right| \\
 \end{array}} {\mbox{z}_\mbox{+} _{ij} q_{ij}^\dag } \mbox{
exp}\sum\limits_{\begin{array}{l}
 \mbox{-}s\leqslant i\leqslant s \\
 \left| i \right|=\left| j \right| \\
 \end{array}} {\left\{ {\lambda _{ij} u_{ij}^ } \right\}}
g_\mbox{0}^{\tfrac{\mbox{N}}{\mbox{2}}} } \right\rangle \\
 \mbox{one can see that } \\
 \frac{\partial ^2H_0 \left( {{\rm {\bf z}}_0 ,{\rm {\bf \bar {z}}}_0 }
\right)}{\partial {\rm {\bf z}}_{+1} \partial {\rm {\bf z}}_{+1}
}=0\Rightarrow \left\langle {g_\mbox{0}^{\tfrac{\mbox{N}}{\mbox{2}}} \left(
{{\rm {\bf z}}_0 } \right)} \mathrel{\left| {\vphantom
{{g_\mbox{0}^{\tfrac{\mbox{N}}{\mbox{2}}} \left( {{\rm {\bf z}}_0 } \right)}
{H_0 q_{i_1 j_1 }^\dag \cdots q_{i_m j_m }^\dag
g_\mbox{0}^{\tfrac{\mbox{N}}{\mbox{2}}} \left( {{\rm {\bf z}}_0 } \right)}}}
\right. \kern-\nulldelimiterspace} {H_0 q_{i_1 j_1 }^\dag \cdots q_{i_m j_m
}^\dag g_\mbox{0}^{\tfrac{\mbox{N}}{\mbox{2}}} \left( {{\rm {\bf z}}_0 }
\right)} \right\rangle =0\mbox{ }\forall m\geqslant 2 \\
 \Rightarrow H_0 \left( {{\rm {\bf z}},{\rm {\bf \bar {z}}}} \right)\mbox{
is first order in }{\rm {\bf z}}\mbox{ and that }H_0 \mbox{ can be expressed
first order in }q_{ij}^\dag \mbox{ } \\
 \mbox{From }\frac{\partial ^2H_0 \left( {{\rm {\bf z}}_0 ,{\rm {\bf \bar
{z}}}_0 } \right)}{\partial {\rm {\bf \bar {z}}}_{+1} \partial {\rm {\bf
\bar {z}}}_{+1} }=0,\mbox{ }\frac{\partial ^2H_0 \left( {{\rm {\bf z}}_0
,{\rm {\bf \bar {z}}}_0 } \right)}{\partial \lambda \partial \lambda
}=0,\mbox{ }\frac{\partial ^2H_0 \left( {{\rm {\bf z}}_0 ,{\rm {\bf \bar
{z}}}_0 } \right)}{\partial \bar {\lambda }\partial \bar {\lambda }}=0\mbox{
one reaches the same } \\
 \mbox{conclusion about the operators }q_{ij}^ \mbox{ }\And \mbox{ }u_{ij}^
.\mbox{ If this were true then } \\
 H_0 \mbox{ =}E_0 \left| {g_\mbox{0}^{\tfrac{\mbox{N}}{\mbox{2}}} }
\right\rangle \left\langle {g_\mbox{0}^{\tfrac{\mbox{N}}{\mbox{2}}} }
\right|+\sum\limits_n {\left( {E_n -E_0 } \right)Q_n^\dag \left|
{g_\mbox{0}^{\tfrac{\mbox{N}}{\mbox{2}}} } \right\rangle \left\langle
{g_\mbox{0}^{\tfrac{\mbox{N}}{\mbox{2}}} } \right|Q_n^ } \\
 \mbox{where }Q_n^\dag =\sum\limits_{\begin{array}{c}
 i-s\leqslant i<j\leqslant s \\
 \left| i \right|\ne \left| j \right| \\
 \end{array}} {Q_{nij}^ q_{ij}^\dag } +\sum\limits_{\begin{array}{c}
 \mbox{-}s\leqslant i\leqslant s \\
 \left| i \right|=\left| j \right| \\
 \end{array}} {U_{nij}^ u_{ij}^ } \mbox{ . } \\
 \end{array}
\]
Thus if the excited states were single excitations out of the reference
state were exact, the off diagonal blocks would be identically zero and the
Hamiltonian could be written as a bilinear form in the above excitation
operators. In this case the quadratic expansion would be exact and Taylor
series quadratic error would be zero. (Note that diagonalizing the Hessian
does not accomplish this unless the off diagonal blocks are actually zero,
as the eigen excitation operators would then be linear combinations of the
excitation operators and their adjoints.) The slower the convergence of the
Taylor series, i.e. the greater the quadratic error term, the poorer the
approximation of excited states by single particle excitations out of the
reference is.

Once the coherent state manifold is defined one is free to use any well
defined coordinate system in which to find the critical points of the energy
function i.e. classical Hamiltonian, for that manifold. We have published a
procedure based on the coordinates based on the following parameterization
of AGP states:
\[
\left| {g^{\frac{\mbox{N}}{\mbox{2}}}} \right\rangle =P_\mbox{N}
e^{\sum\limits_{1\leqslant i\leqslant s=\raise0.7ex\hbox{$r$}
\!\mathord{\left/ {\vphantom {r
2}}\right.\kern-\nulldelimiterspace}\!\lower0.7ex\hbox{$2$}} {g_{ij}
\mbox{a}_i^\dag \mbox{a}_j^\dag } }\left| \phi \right\rangle \mbox{ }
\]
This procedure will be compared to ones based on coordinate system
{\{}\textbf{z}{\}}.

\underline {\textbf{5. Functions defined on the Coset Space. }}

\underline {\textbf{5.1 Space of Holomorphic Quantum Wavefunctions}}

Due to the irreducibility of the group G on the space $H$ this set of states
forms an overcomplete basis in $H$, thus one has the resolution of the
identity
\[
{\rm {\bf I}}=\int {\left| {\rm {\bf z}} \right\rangle \left\langle {\rm
{\bf z}} \right|} d\mu \left( {\rm {\bf z}} \right)
\]
for a suitably chosen measure $\mu \left( {\rm {\bf z}} \right)$. Using this
resolution we obtain
\[
\mbox{ }\left| \Psi \right\rangle =\int {\left| {\rm {\bf z}} \right\rangle
} \left\langle {{\rm {\bf z}}} \mathrel{\left| {\vphantom {{\rm {\bf z}}
\Psi }} \right. \kern-\nulldelimiterspace} {\Psi } \right\rangle d\mu \left(
{\rm {\bf z}} \right)=\int {\mbox{ }\left| {\rm {\bf z}} \right\rangle }
\Psi \left( {\rm {\bf z}} \right)d\mu \left( {\rm {\bf z}} \right)
\]
where the amplitudes $\Psi \left( {\rm {\bf z}} \right)$ are holomorphic and
are in 1-1 correspondence with the states $\left| \Psi \right\rangle $. The
space, $H\left( M \right)$, of all such amplitudes is a Hilbert space that
is isomorphic to $H$. Thus all quantum mechanical operators that act in $H$
have integral kernel representations in the space of holomorphic amplitudes
given by
\[
\left\langle {\Psi _1 } \mathrel{\left| {\vphantom {{\Psi _1 } {\mbox{H}\Psi
_2 }}} \right. \kern-\nulldelimiterspace} {\mbox{H}\Psi _2 } \right\rangle
=\int {\mbox{ }\bar {\Psi }_1 \left( {{\rm {\bf {z}'}}} \right)}
\mbox{H}\left( {{\rm {\bf {z}'}},{\rm {\bf z}}} \right)\Psi _2 \left( {\rm
{\bf z}} \right)d\mu \left( {{\rm {\bf {z}'}}} \right)d\mu \left( {\rm {\bf
z}} \right)
\]
\[
\mbox{H}\left( {{\rm {\bf {z}'}},{\rm {\bf z}}} \right)=\left\langle {{\rm
{\bf {z}'}}} \mathrel{\left| {\vphantom {{{\rm {\bf {z}'}}} {\mbox{H}{\rm
{\bf z}}}}} \right. \kern-\nulldelimiterspace} {\mbox{H}{\rm {\bf z}}}
\right\rangle \mbox{ }
\]
$\mbox{H}\left( {{\rm {\bf {z}'}},{\rm {\bf z}}} \right)\mbox{ can be
constructed as the analytic extension of}
\mbox{H}\left( {{\rm {\bf z}},{\rm {\bf z}}} \right),$ which is the
unnormalized energy functional used in the classical equations of motion.

\underline {\textbf{5.2 Representation of Observables as Differential
Operators. }}

Elements of the universal enveloping algebra of $sl\left( {H\mbox{ }}
\right)$ act as differential operators on functions defined on $M$. As
quantum mechanical observables formally belong to the enveloping algebra,
they too can be expressed as differential operators, and their expectation
values as the value of the action of these operators at points in $M$. In
particular
\[
\left\langle {g^{\frac{\mbox{N}}{\mbox{2}}}\left( {\rm {\bf z}} \right)}
\mathrel{\left| {\vphantom {{g^{\frac{\mbox{N}}{\mbox{2}}}\left( {\rm {\bf
z}} \right)} {\mbox{a}_i^\dag \mbox{a}_j g^{\frac{\mbox{N}}{\mbox{2}}}\left(
{\rm {\bf z}} \right)}}} \right. \kern-\nulldelimiterspace} {\mbox{a}_i^\dag
\mbox{a}_j g^{\frac{\mbox{N}}{\mbox{2}}}\left( {\rm {\bf z}} \right)}
\right\rangle =P_{ij} \left( {\frac{\partial }{\partial {\rm {\bf
z}}},\frac{\partial }{\partial {\rm {\bf \bar {z}}}}} \right)\left\langle
{g^{\frac{\mbox{N}}{\mbox{2}}}\left( {\rm {\bf z}} \right)} \mathrel{\left|
{\vphantom {{g^{\frac{\mbox{N}}{\mbox{2}}}\left( {\rm {\bf z}} \right)}
{g^{\frac{\mbox{N}}{\mbox{2}}}\left( {\rm {\bf z}} \right)}}} \right.
\kern-\nulldelimiterspace} {g^{\frac{\mbox{N}}{\mbox{2}}}\left( {\rm {\bf
z}} \right)} \right\rangle
\]
$\mbox{where }P_{ij} \left( {\frac{\partial }{\partial {\rm {\bf
z}}},\frac{\partial }{\partial {\rm {\bf \bar {z}}}}} \right)\mbox{ is a
polynomial in }\left\{ {\frac{\partial }{\partial {\rm {\bf
z}}},\frac{\partial }{\partial {\rm {\bf \bar {z}}}}} \right\}.$ Such
relationships are important in constructing efficient expressions for
expectation values, as well as for the various terms appearing the in
generalized classical equations of motion on the coherent state manifold and
the associated linear response functions.

Thus on the space $H\left( M \right)$ they can be represented both by
integral kernels as well as differential operators. There is and has been an
active research effort in classifying and analyzing the properties of
differential operators on coset manifolds in an abstract setting. We propose
to apply this work to the study of quantum mechanical observables involved
in response functions obtained from 2 particle propagators.

\underline {\textbf{5.3 Induced Representations}}

Elements of G act on the coset space mapping it to itself, thus give rise to
a, in general \textit{nonlinear, }representation of G. Thus
\[
\begin{array}{c}
 G:\mbox{ }\raise0.7ex\hbox{$G$} \!\mathord{\left/ {\vphantom {G
H}}\right.\kern-\nulldelimiterspace}\!\lower0.7ex\hbox{$H$}\to
\raise0.7ex\hbox{$G$} \!\mathord{\left/ {\vphantom {G
H}}\right.\kern-\nulldelimiterspace}\!\lower0.7ex\hbox{$H$}\mbox{
}\Rightarrow \mbox{ }\sigma \left( g \right):M\to M\mbox{ i.e. }\sigma
\left( g \right){\rm {\bf z}}={\rm {\bf {z}'}} \\
 \sigma \left( h \right),\mbox{ }h\in H\mbox{ leaves each point in }M\mbox{
invariant} \\
 \mbox{Thus a new representation of G is induced on }H\left( M \right)\mbox{
by} \\
 \Pi \left( \mbox{g} \right)\Psi \left( {\rm {\bf z}} \right)=\Psi \left(
{\sigma \left( g \right){\rm {\bf z}}} \right) \\
 \end{array}
\]
The induced representation depends on the original representation
$\mbox{G}_H $ and the isotropy subgroup $\mbox{H}_V $ , which itself
corresponds to the form of the coherent state. All but two of the classical
Riemannian symmetric coset spaces can be expressed as exponential of off
diagonal matrices, for these cases the action of the group on the coset
space is represented as fractional linear transformations
\[
\sigma \left( g \right)\mbox{ }{\rm {\bf z}}=\left( {{\rm {\bf Az}}+{\rm
{\bf B}}} \right)\left( {{\rm {\bf Cz}}+{\rm {\bf D}}} \right)^{-1}={\rm
{\bf {z}'}}
\]
where $\left( {{\begin{array}{*{20}c}
 A \hfill & B \hfill \\
 C \hfill & D \hfill \\
\end{array} }} \right)$ is the defining representation of $g\mbox{ }\in G$.
However for the other two cosets $\raise0.7ex\hbox{${\mbox{SU(r)}}$}
\!\mathord{\left/ {\vphantom {{\mbox{SU(r)}}
{\mbox{USp(2s)}}}}\right.\kern-\nulldelimiterspace}\!\lower0.7ex\hbox{${\mbox{USp(2s)}}$}$
and $\raise0.7ex\hbox{${\mbox{SU(r)}}$} \!\mathord{\left/ {\vphantom
{{\mbox{SU(r)}}
{\mbox{SO(r)}}}}\right.\kern-\nulldelimiterspace}\!\lower0.7ex\hbox{${\mbox{SO(r)}}$}$
this is not the case. The coset $\raise0.7ex\hbox{${\mbox{SU(r)}}$}
\!\mathord{\left/ {\vphantom {{\mbox{SU(r)}}
{\mbox{USp(2s)}}}}\right.\kern-\nulldelimiterspace}\!\lower0.7ex\hbox{${\mbox{USp(2s)}}$}$
is the primary focus of this study, which will investigate the details of
the action of $\mbox{SU(r) on }\raise0.7ex\hbox{${\mbox{SU(r)}}$}
\!\mathord{\left/ {\vphantom {{\mbox{SU(r)}}
{\mbox{USp(2s)}}}}\right.\kern-\nulldelimiterspace}\!\lower0.7ex\hbox{${\mbox{USp(2s)}}$}\mbox{
and SL(r) on its complexification }\raise0.7ex\hbox{${\mbox{SL(r)}}$}
\!\mathord{\left/ {\vphantom {{\mbox{SL(r)}}
{\mbox{Sp(2s,C)}}}}\right.\kern-\nulldelimiterspace}\!\lower0.7ex\hbox{${\mbox{Sp(2s,C)}}$}$.
The action of the group $\mbox{SU(r)}$ on the coset
$\raise0.7ex\hbox{${\mbox{SU(r)}}$} \!\mathord{\left/ {\vphantom
{{\mbox{SU(r)}}
{\mbox{USp(2s)}}}}\right.\kern-\nulldelimiterspace}\!\lower0.7ex\hbox{${\mbox{USp(2s)}}$}$
is important in the study of physical symmetries, such as spin and angular
momentum. We have shown how one can produce pure symmetry coherent states in
the case of $\raise0.7ex\hbox{${\mbox{SU(r)}}$} \!\mathord{\left/ {\vphantom
{{\mbox{SU(r)}} {\mbox{U(N)}\otimes
\mbox{U(r-N)}}}}\right.\kern-\nulldelimiterspace}\!\lower0.7ex\hbox{${\mbox{U(N)}\otimes
\mbox{U(r-N)}}$}$ coherent states and in this work will extend such a
construction to $\raise0.7ex\hbox{${\mbox{SU(r)}}$} \!\mathord{\left/
{\vphantom {{\mbox{SU(r)}}
{\mbox{USp(2s)}}}}\right.\kern-\nulldelimiterspace}\!\lower0.7ex\hbox{${\mbox{USp(2s)}}$}$
ones.

\end{document}
